\documentclass[11pt]{article}

    \usepackage[breakable]{tcolorbox}
    \usepackage{parskip} % Stop auto-indenting (to mimic markdown behaviour)
    

    % Basic figure setup, for now with no caption control since it's done
    % automatically by Pandoc (which extracts ![](path) syntax from Markdown).
    \usepackage{graphicx}
    % Keep aspect ratio if custom image width or height is specified
    \setkeys{Gin}{keepaspectratio}
    % Maintain compatibility with old templates. Remove in nbconvert 6.0
    \let\Oldincludegraphics\includegraphics
    % Ensure that by default, figures have no caption (until we provide a
    % proper Figure object with a Caption API and a way to capture that
    % in the conversion process - todo).
    \usepackage{caption}
    \DeclareCaptionFormat{nocaption}{}
    \captionsetup{format=nocaption,aboveskip=0pt,belowskip=0pt}

    \usepackage{float}
    \floatplacement{figure}{H} % forces figures to be placed at the correct location
    \usepackage{xcolor} % Allow colors to be defined
    \usepackage{enumerate} % Needed for markdown enumerations to work
    \usepackage{geometry} % Used to adjust the document margins
    \usepackage{amsmath} % Equations
    \usepackage{amssymb} % Equations
    \usepackage{textcomp} % defines textquotesingle
    % Hack from http://tex.stackexchange.com/a/47451/13684:
    \AtBeginDocument{%
        \def\PYZsq{\textquotesingle}% Upright quotes in Pygmentized code
    }
    \usepackage{upquote} % Upright quotes for verbatim code
    \usepackage{eurosym} % defines \euro

    \usepackage{iftex}
    \ifPDFTeX
        \usepackage[T1]{fontenc}
        \IfFileExists{alphabeta.sty}{
              \usepackage{alphabeta}
          }{
              \usepackage[mathletters]{ucs}
              \usepackage[utf8x]{inputenc}
          }
    \else
        \usepackage{fontspec}
        \usepackage{unicode-math}
    \fi

    \usepackage{fancyvrb} % verbatim replacement that allows latex
    \usepackage{grffile} % extends the file name processing of package graphics
                         % to support a larger range
    \makeatletter % fix for old versions of grffile with XeLaTeX
    \@ifpackagelater{grffile}{2019/11/01}
    {
      % Do nothing on new versions
    }
    {
      \def\Gread@@xetex#1{%
        \IfFileExists{"\Gin@base".bb}%
        {\Gread@eps{\Gin@base.bb}}%
        {\Gread@@xetex@aux#1}%
      }
    }
    \makeatother
    \usepackage[Export]{adjustbox} % Used to constrain images to a maximum size
    \adjustboxset{max size={0.9\linewidth}{0.9\paperheight}}

    % The hyperref package gives us a pdf with properly built
    % internal navigation ('pdf bookmarks' for the table of contents,
    % internal cross-reference links, web links for URLs, etc.)
    \usepackage{hyperref}
    % The default LaTeX title has an obnoxious amount of whitespace. By default,
    % titling removes some of it. It also provides customization options.
    \usepackage{titling}
    \usepackage{longtable} % longtable support required by pandoc >1.10
    \usepackage{booktabs}  % table support for pandoc > 1.12.2
    \usepackage{array}     % table support for pandoc >= 2.11.3
    \usepackage{calc}      % table minipage width calculation for pandoc >= 2.11.1
    \usepackage[inline]{enumitem} % IRkernel/repr support (it uses the enumerate* environment)
    \usepackage[normalem]{ulem} % ulem is needed to support strikethroughs (\sout)
                                % normalem makes italics be italics, not underlines
    \usepackage{soul}      % strikethrough (\st) support for pandoc >= 3.0.0
    \usepackage{mathrsfs}
    

    
    % Colors for the hyperref package
    \definecolor{urlcolor}{rgb}{0,.145,.698}
    \definecolor{linkcolor}{rgb}{.71,0.21,0.01}
    \definecolor{citecolor}{rgb}{.12,.54,.11}

    % ANSI colors
    \definecolor{ansi-black}{HTML}{3E424D}
    \definecolor{ansi-black-intense}{HTML}{282C36}
    \definecolor{ansi-red}{HTML}{E75C58}
    \definecolor{ansi-red-intense}{HTML}{B22B31}
    \definecolor{ansi-green}{HTML}{00A250}
    \definecolor{ansi-green-intense}{HTML}{007427}
    \definecolor{ansi-yellow}{HTML}{DDB62B}
    \definecolor{ansi-yellow-intense}{HTML}{B27D12}
    \definecolor{ansi-blue}{HTML}{208FFB}
    \definecolor{ansi-blue-intense}{HTML}{0065CA}
    \definecolor{ansi-magenta}{HTML}{D160C4}
    \definecolor{ansi-magenta-intense}{HTML}{A03196}
    \definecolor{ansi-cyan}{HTML}{60C6C8}
    \definecolor{ansi-cyan-intense}{HTML}{258F8F}
    \definecolor{ansi-white}{HTML}{C5C1B4}
    \definecolor{ansi-white-intense}{HTML}{A1A6B2}
    \definecolor{ansi-default-inverse-fg}{HTML}{FFFFFF}
    \definecolor{ansi-default-inverse-bg}{HTML}{000000}

    % common color for the border for error outputs.
    \definecolor{outerrorbackground}{HTML}{FFDFDF}

    % commands and environments needed by pandoc snippets
    % extracted from the output of `pandoc -s`
    \providecommand{\tightlist}{%
      \setlength{\itemsep}{0pt}\setlength{\parskip}{0pt}}
    \DefineVerbatimEnvironment{Highlighting}{Verbatim}{commandchars=\\\{\}}
    % Add ',fontsize=\small' for more characters per line
    \newenvironment{Shaded}{}{}
    \newcommand{\KeywordTok}[1]{\textcolor[rgb]{0.00,0.44,0.13}{\textbf{{#1}}}}
    \newcommand{\DataTypeTok}[1]{\textcolor[rgb]{0.56,0.13,0.00}{{#1}}}
    \newcommand{\DecValTok}[1]{\textcolor[rgb]{0.25,0.63,0.44}{{#1}}}
    \newcommand{\BaseNTok}[1]{\textcolor[rgb]{0.25,0.63,0.44}{{#1}}}
    \newcommand{\FloatTok}[1]{\textcolor[rgb]{0.25,0.63,0.44}{{#1}}}
    \newcommand{\CharTok}[1]{\textcolor[rgb]{0.25,0.44,0.63}{{#1}}}
    \newcommand{\StringTok}[1]{\textcolor[rgb]{0.25,0.44,0.63}{{#1}}}
    \newcommand{\CommentTok}[1]{\textcolor[rgb]{0.38,0.63,0.69}{\textit{{#1}}}}
    \newcommand{\OtherTok}[1]{\textcolor[rgb]{0.00,0.44,0.13}{{#1}}}
    \newcommand{\AlertTok}[1]{\textcolor[rgb]{1.00,0.00,0.00}{\textbf{{#1}}}}
    \newcommand{\FunctionTok}[1]{\textcolor[rgb]{0.02,0.16,0.49}{{#1}}}
    \newcommand{\RegionMarkerTok}[1]{{#1}}
    \newcommand{\ErrorTok}[1]{\textcolor[rgb]{1.00,0.00,0.00}{\textbf{{#1}}}}
    \newcommand{\NormalTok}[1]{{#1}}

    % Additional commands for more recent versions of Pandoc
    \newcommand{\ConstantTok}[1]{\textcolor[rgb]{0.53,0.00,0.00}{{#1}}}
    \newcommand{\SpecialCharTok}[1]{\textcolor[rgb]{0.25,0.44,0.63}{{#1}}}
    \newcommand{\VerbatimStringTok}[1]{\textcolor[rgb]{0.25,0.44,0.63}{{#1}}}
    \newcommand{\SpecialStringTok}[1]{\textcolor[rgb]{0.73,0.40,0.53}{{#1}}}
    \newcommand{\ImportTok}[1]{{#1}}
    \newcommand{\DocumentationTok}[1]{\textcolor[rgb]{0.73,0.13,0.13}{\textit{{#1}}}}
    \newcommand{\AnnotationTok}[1]{\textcolor[rgb]{0.38,0.63,0.69}{\textbf{\textit{{#1}}}}}
    \newcommand{\CommentVarTok}[1]{\textcolor[rgb]{0.38,0.63,0.69}{\textbf{\textit{{#1}}}}}
    \newcommand{\VariableTok}[1]{\textcolor[rgb]{0.10,0.09,0.49}{{#1}}}
    \newcommand{\ControlFlowTok}[1]{\textcolor[rgb]{0.00,0.44,0.13}{\textbf{{#1}}}}
    \newcommand{\OperatorTok}[1]{\textcolor[rgb]{0.40,0.40,0.40}{{#1}}}
    \newcommand{\BuiltInTok}[1]{{#1}}
    \newcommand{\ExtensionTok}[1]{{#1}}
    \newcommand{\PreprocessorTok}[1]{\textcolor[rgb]{0.74,0.48,0.00}{{#1}}}
    \newcommand{\AttributeTok}[1]{\textcolor[rgb]{0.49,0.56,0.16}{{#1}}}
    \newcommand{\InformationTok}[1]{\textcolor[rgb]{0.38,0.63,0.69}{\textbf{\textit{{#1}}}}}
    \newcommand{\WarningTok}[1]{\textcolor[rgb]{0.38,0.63,0.69}{\textbf{\textit{{#1}}}}}
    \makeatletter
    \newsavebox\pandoc@box
    \newcommand*\pandocbounded[1]{%
      \sbox\pandoc@box{#1}%
      % scaling factors for width and height
      \Gscale@div\@tempa\textheight{\dimexpr\ht\pandoc@box+\dp\pandoc@box\relax}%
      \Gscale@div\@tempb\linewidth{\wd\pandoc@box}%
      % select the smaller of both
      \ifdim\@tempb\p@<\@tempa\p@
        \let\@tempa\@tempb
      \fi
      % scaling accordingly (\@tempa < 1)
      \ifdim\@tempa\p@<\p@
        \scalebox{\@tempa}{\usebox\pandoc@box}%
      % scaling not needed, use as it is
      \else
        \usebox{\pandoc@box}%
      \fi
    }
    \makeatother

    % Define a nice break command that doesn't care if a line doesn't already
    % exist.
    \def\br{\hspace*{\fill} \\* }
    % Math Jax compatibility definitions
    \def\gt{>}
    \def\lt{<}
    \let\Oldtex\TeX
    \let\Oldlatex\LaTeX
    \renewcommand{\TeX}{\textrm{\Oldtex}}
    \renewcommand{\LaTeX}{\textrm{\Oldlatex}}
    % Document parameters
    % Document title
    \title{exploring\_word\_vectors}
    
    
    
    
    
    
    
% Pygments definitions
\makeatletter
\def\PY@reset{\let\PY@it=\relax \let\PY@bf=\relax%
    \let\PY@ul=\relax \let\PY@tc=\relax%
    \let\PY@bc=\relax \let\PY@ff=\relax}
\def\PY@tok#1{\csname PY@tok@#1\endcsname}
\def\PY@toks#1+{\ifx\relax#1\empty\else%
    \PY@tok{#1}\expandafter\PY@toks\fi}
\def\PY@do#1{\PY@bc{\PY@tc{\PY@ul{%
    \PY@it{\PY@bf{\PY@ff{#1}}}}}}}
\def\PY#1#2{\PY@reset\PY@toks#1+\relax+\PY@do{#2}}

\@namedef{PY@tok@w}{\def\PY@tc##1{\textcolor[rgb]{0.73,0.73,0.73}{##1}}}
\@namedef{PY@tok@c}{\let\PY@it=\textit\def\PY@tc##1{\textcolor[rgb]{0.24,0.48,0.48}{##1}}}
\@namedef{PY@tok@cp}{\def\PY@tc##1{\textcolor[rgb]{0.61,0.40,0.00}{##1}}}
\@namedef{PY@tok@k}{\let\PY@bf=\textbf\def\PY@tc##1{\textcolor[rgb]{0.00,0.50,0.00}{##1}}}
\@namedef{PY@tok@kp}{\def\PY@tc##1{\textcolor[rgb]{0.00,0.50,0.00}{##1}}}
\@namedef{PY@tok@kt}{\def\PY@tc##1{\textcolor[rgb]{0.69,0.00,0.25}{##1}}}
\@namedef{PY@tok@o}{\def\PY@tc##1{\textcolor[rgb]{0.40,0.40,0.40}{##1}}}
\@namedef{PY@tok@ow}{\let\PY@bf=\textbf\def\PY@tc##1{\textcolor[rgb]{0.67,0.13,1.00}{##1}}}
\@namedef{PY@tok@nb}{\def\PY@tc##1{\textcolor[rgb]{0.00,0.50,0.00}{##1}}}
\@namedef{PY@tok@nf}{\def\PY@tc##1{\textcolor[rgb]{0.00,0.00,1.00}{##1}}}
\@namedef{PY@tok@nc}{\let\PY@bf=\textbf\def\PY@tc##1{\textcolor[rgb]{0.00,0.00,1.00}{##1}}}
\@namedef{PY@tok@nn}{\let\PY@bf=\textbf\def\PY@tc##1{\textcolor[rgb]{0.00,0.00,1.00}{##1}}}
\@namedef{PY@tok@ne}{\let\PY@bf=\textbf\def\PY@tc##1{\textcolor[rgb]{0.80,0.25,0.22}{##1}}}
\@namedef{PY@tok@nv}{\def\PY@tc##1{\textcolor[rgb]{0.10,0.09,0.49}{##1}}}
\@namedef{PY@tok@no}{\def\PY@tc##1{\textcolor[rgb]{0.53,0.00,0.00}{##1}}}
\@namedef{PY@tok@nl}{\def\PY@tc##1{\textcolor[rgb]{0.46,0.46,0.00}{##1}}}
\@namedef{PY@tok@ni}{\let\PY@bf=\textbf\def\PY@tc##1{\textcolor[rgb]{0.44,0.44,0.44}{##1}}}
\@namedef{PY@tok@na}{\def\PY@tc##1{\textcolor[rgb]{0.41,0.47,0.13}{##1}}}
\@namedef{PY@tok@nt}{\let\PY@bf=\textbf\def\PY@tc##1{\textcolor[rgb]{0.00,0.50,0.00}{##1}}}
\@namedef{PY@tok@nd}{\def\PY@tc##1{\textcolor[rgb]{0.67,0.13,1.00}{##1}}}
\@namedef{PY@tok@s}{\def\PY@tc##1{\textcolor[rgb]{0.73,0.13,0.13}{##1}}}
\@namedef{PY@tok@sd}{\let\PY@it=\textit\def\PY@tc##1{\textcolor[rgb]{0.73,0.13,0.13}{##1}}}
\@namedef{PY@tok@si}{\let\PY@bf=\textbf\def\PY@tc##1{\textcolor[rgb]{0.64,0.35,0.47}{##1}}}
\@namedef{PY@tok@se}{\let\PY@bf=\textbf\def\PY@tc##1{\textcolor[rgb]{0.67,0.36,0.12}{##1}}}
\@namedef{PY@tok@sr}{\def\PY@tc##1{\textcolor[rgb]{0.64,0.35,0.47}{##1}}}
\@namedef{PY@tok@ss}{\def\PY@tc##1{\textcolor[rgb]{0.10,0.09,0.49}{##1}}}
\@namedef{PY@tok@sx}{\def\PY@tc##1{\textcolor[rgb]{0.00,0.50,0.00}{##1}}}
\@namedef{PY@tok@m}{\def\PY@tc##1{\textcolor[rgb]{0.40,0.40,0.40}{##1}}}
\@namedef{PY@tok@gh}{\let\PY@bf=\textbf\def\PY@tc##1{\textcolor[rgb]{0.00,0.00,0.50}{##1}}}
\@namedef{PY@tok@gu}{\let\PY@bf=\textbf\def\PY@tc##1{\textcolor[rgb]{0.50,0.00,0.50}{##1}}}
\@namedef{PY@tok@gd}{\def\PY@tc##1{\textcolor[rgb]{0.63,0.00,0.00}{##1}}}
\@namedef{PY@tok@gi}{\def\PY@tc##1{\textcolor[rgb]{0.00,0.52,0.00}{##1}}}
\@namedef{PY@tok@gr}{\def\PY@tc##1{\textcolor[rgb]{0.89,0.00,0.00}{##1}}}
\@namedef{PY@tok@ge}{\let\PY@it=\textit}
\@namedef{PY@tok@gs}{\let\PY@bf=\textbf}
\@namedef{PY@tok@ges}{\let\PY@bf=\textbf\let\PY@it=\textit}
\@namedef{PY@tok@gp}{\let\PY@bf=\textbf\def\PY@tc##1{\textcolor[rgb]{0.00,0.00,0.50}{##1}}}
\@namedef{PY@tok@go}{\def\PY@tc##1{\textcolor[rgb]{0.44,0.44,0.44}{##1}}}
\@namedef{PY@tok@gt}{\def\PY@tc##1{\textcolor[rgb]{0.00,0.27,0.87}{##1}}}
\@namedef{PY@tok@err}{\def\PY@bc##1{{\setlength{\fboxsep}{\string -\fboxrule}\fcolorbox[rgb]{1.00,0.00,0.00}{1,1,1}{\strut ##1}}}}
\@namedef{PY@tok@kc}{\let\PY@bf=\textbf\def\PY@tc##1{\textcolor[rgb]{0.00,0.50,0.00}{##1}}}
\@namedef{PY@tok@kd}{\let\PY@bf=\textbf\def\PY@tc##1{\textcolor[rgb]{0.00,0.50,0.00}{##1}}}
\@namedef{PY@tok@kn}{\let\PY@bf=\textbf\def\PY@tc##1{\textcolor[rgb]{0.00,0.50,0.00}{##1}}}
\@namedef{PY@tok@kr}{\let\PY@bf=\textbf\def\PY@tc##1{\textcolor[rgb]{0.00,0.50,0.00}{##1}}}
\@namedef{PY@tok@bp}{\def\PY@tc##1{\textcolor[rgb]{0.00,0.50,0.00}{##1}}}
\@namedef{PY@tok@fm}{\def\PY@tc##1{\textcolor[rgb]{0.00,0.00,1.00}{##1}}}
\@namedef{PY@tok@vc}{\def\PY@tc##1{\textcolor[rgb]{0.10,0.09,0.49}{##1}}}
\@namedef{PY@tok@vg}{\def\PY@tc##1{\textcolor[rgb]{0.10,0.09,0.49}{##1}}}
\@namedef{PY@tok@vi}{\def\PY@tc##1{\textcolor[rgb]{0.10,0.09,0.49}{##1}}}
\@namedef{PY@tok@vm}{\def\PY@tc##1{\textcolor[rgb]{0.10,0.09,0.49}{##1}}}
\@namedef{PY@tok@sa}{\def\PY@tc##1{\textcolor[rgb]{0.73,0.13,0.13}{##1}}}
\@namedef{PY@tok@sb}{\def\PY@tc##1{\textcolor[rgb]{0.73,0.13,0.13}{##1}}}
\@namedef{PY@tok@sc}{\def\PY@tc##1{\textcolor[rgb]{0.73,0.13,0.13}{##1}}}
\@namedef{PY@tok@dl}{\def\PY@tc##1{\textcolor[rgb]{0.73,0.13,0.13}{##1}}}
\@namedef{PY@tok@s2}{\def\PY@tc##1{\textcolor[rgb]{0.73,0.13,0.13}{##1}}}
\@namedef{PY@tok@sh}{\def\PY@tc##1{\textcolor[rgb]{0.73,0.13,0.13}{##1}}}
\@namedef{PY@tok@s1}{\def\PY@tc##1{\textcolor[rgb]{0.73,0.13,0.13}{##1}}}
\@namedef{PY@tok@mb}{\def\PY@tc##1{\textcolor[rgb]{0.40,0.40,0.40}{##1}}}
\@namedef{PY@tok@mf}{\def\PY@tc##1{\textcolor[rgb]{0.40,0.40,0.40}{##1}}}
\@namedef{PY@tok@mh}{\def\PY@tc##1{\textcolor[rgb]{0.40,0.40,0.40}{##1}}}
\@namedef{PY@tok@mi}{\def\PY@tc##1{\textcolor[rgb]{0.40,0.40,0.40}{##1}}}
\@namedef{PY@tok@il}{\def\PY@tc##1{\textcolor[rgb]{0.40,0.40,0.40}{##1}}}
\@namedef{PY@tok@mo}{\def\PY@tc##1{\textcolor[rgb]{0.40,0.40,0.40}{##1}}}
\@namedef{PY@tok@ch}{\let\PY@it=\textit\def\PY@tc##1{\textcolor[rgb]{0.24,0.48,0.48}{##1}}}
\@namedef{PY@tok@cm}{\let\PY@it=\textit\def\PY@tc##1{\textcolor[rgb]{0.24,0.48,0.48}{##1}}}
\@namedef{PY@tok@cpf}{\let\PY@it=\textit\def\PY@tc##1{\textcolor[rgb]{0.24,0.48,0.48}{##1}}}
\@namedef{PY@tok@c1}{\let\PY@it=\textit\def\PY@tc##1{\textcolor[rgb]{0.24,0.48,0.48}{##1}}}
\@namedef{PY@tok@cs}{\let\PY@it=\textit\def\PY@tc##1{\textcolor[rgb]{0.24,0.48,0.48}{##1}}}

\def\PYZbs{\char`\\}
\def\PYZus{\char`\_}
\def\PYZob{\char`\{}
\def\PYZcb{\char`\}}
\def\PYZca{\char`\^}
\def\PYZam{\char`\&}
\def\PYZlt{\char`\<}
\def\PYZgt{\char`\>}
\def\PYZsh{\char`\#}
\def\PYZpc{\char`\%}
\def\PYZdl{\char`\$}
\def\PYZhy{\char`\-}
\def\PYZsq{\char`\'}
\def\PYZdq{\char`\"}
\def\PYZti{\char`\~}
% for compatibility with earlier versions
\def\PYZat{@}
\def\PYZlb{[}
\def\PYZrb{]}
\makeatother


    % For linebreaks inside Verbatim environment from package fancyvrb.
    \makeatletter
        \newbox\Wrappedcontinuationbox
        \newbox\Wrappedvisiblespacebox
        \newcommand*\Wrappedvisiblespace {\textcolor{red}{\textvisiblespace}}
        \newcommand*\Wrappedcontinuationsymbol {\textcolor{red}{\llap{\tiny$\m@th\hookrightarrow$}}}
        \newcommand*\Wrappedcontinuationindent {3ex }
        \newcommand*\Wrappedafterbreak {\kern\Wrappedcontinuationindent\copy\Wrappedcontinuationbox}
        % Take advantage of the already applied Pygments mark-up to insert
        % potential linebreaks for TeX processing.
        %        {, <, #, %, $, ' and ": go to next line.
        %        _, }, ^, &, >, - and ~: stay at end of broken line.
        % Use of \textquotesingle for straight quote.
        \newcommand*\Wrappedbreaksatspecials {%
            \def\PYGZus{\discretionary{\char`\_}{\Wrappedafterbreak}{\char`\_}}%
            \def\PYGZob{\discretionary{}{\Wrappedafterbreak\char`\{}{\char`\{}}%
            \def\PYGZcb{\discretionary{\char`\}}{\Wrappedafterbreak}{\char`\}}}%
            \def\PYGZca{\discretionary{\char`\^}{\Wrappedafterbreak}{\char`\^}}%
            \def\PYGZam{\discretionary{\char`\&}{\Wrappedafterbreak}{\char`\&}}%
            \def\PYGZlt{\discretionary{}{\Wrappedafterbreak\char`\<}{\char`\<}}%
            \def\PYGZgt{\discretionary{\char`\>}{\Wrappedafterbreak}{\char`\>}}%
            \def\PYGZsh{\discretionary{}{\Wrappedafterbreak\char`\#}{\char`\#}}%
            \def\PYGZpc{\discretionary{}{\Wrappedafterbreak\char`\%}{\char`\%}}%
            \def\PYGZdl{\discretionary{}{\Wrappedafterbreak\char`\$}{\char`\$}}%
            \def\PYGZhy{\discretionary{\char`\-}{\Wrappedafterbreak}{\char`\-}}%
            \def\PYGZsq{\discretionary{}{\Wrappedafterbreak\textquotesingle}{\textquotesingle}}%
            \def\PYGZdq{\discretionary{}{\Wrappedafterbreak\char`\"}{\char`\"}}%
            \def\PYGZti{\discretionary{\char`\~}{\Wrappedafterbreak}{\char`\~}}%
        }
        % Some characters . , ; ? ! / are not pygmentized.
        % This macro makes them "active" and they will insert potential linebreaks
        \newcommand*\Wrappedbreaksatpunct {%
            \lccode`\~`\.\lowercase{\def~}{\discretionary{\hbox{\char`\.}}{\Wrappedafterbreak}{\hbox{\char`\.}}}%
            \lccode`\~`\,\lowercase{\def~}{\discretionary{\hbox{\char`\,}}{\Wrappedafterbreak}{\hbox{\char`\,}}}%
            \lccode`\~`\;\lowercase{\def~}{\discretionary{\hbox{\char`\;}}{\Wrappedafterbreak}{\hbox{\char`\;}}}%
            \lccode`\~`\:\lowercase{\def~}{\discretionary{\hbox{\char`\:}}{\Wrappedafterbreak}{\hbox{\char`\:}}}%
            \lccode`\~`\?\lowercase{\def~}{\discretionary{\hbox{\char`\?}}{\Wrappedafterbreak}{\hbox{\char`\?}}}%
            \lccode`\~`\!\lowercase{\def~}{\discretionary{\hbox{\char`\!}}{\Wrappedafterbreak}{\hbox{\char`\!}}}%
            \lccode`\~`\/\lowercase{\def~}{\discretionary{\hbox{\char`\/}}{\Wrappedafterbreak}{\hbox{\char`\/}}}%
            \catcode`\.\active
            \catcode`\,\active
            \catcode`\;\active
            \catcode`\:\active
            \catcode`\?\active
            \catcode`\!\active
            \catcode`\/\active
            \lccode`\~`\~
        }
    \makeatother

    \let\OriginalVerbatim=\Verbatim
    \makeatletter
    \renewcommand{\Verbatim}[1][1]{%
        %\parskip\z@skip
        \sbox\Wrappedcontinuationbox {\Wrappedcontinuationsymbol}%
        \sbox\Wrappedvisiblespacebox {\FV@SetupFont\Wrappedvisiblespace}%
        \def\FancyVerbFormatLine ##1{\hsize\linewidth
            \vtop{\raggedright\hyphenpenalty\z@\exhyphenpenalty\z@
                \doublehyphendemerits\z@\finalhyphendemerits\z@
                \strut ##1\strut}%
        }%
        % If the linebreak is at a space, the latter will be displayed as visible
        % space at end of first line, and a continuation symbol starts next line.
        % Stretch/shrink are however usually zero for typewriter font.
        \def\FV@Space {%
            \nobreak\hskip\z@ plus\fontdimen3\font minus\fontdimen4\font
            \discretionary{\copy\Wrappedvisiblespacebox}{\Wrappedafterbreak}
            {\kern\fontdimen2\font}%
        }%

        % Allow breaks at special characters using \PYG... macros.
        \Wrappedbreaksatspecials
        % Breaks at punctuation characters . , ; ? ! and / need catcode=\active
        \OriginalVerbatim[#1,codes*=\Wrappedbreaksatpunct]%
    }
    \makeatother

    % Exact colors from NB
    \definecolor{incolor}{HTML}{303F9F}
    \definecolor{outcolor}{HTML}{D84315}
    \definecolor{cellborder}{HTML}{CFCFCF}
    \definecolor{cellbackground}{HTML}{F7F7F7}

    % prompt
    \makeatletter
    \newcommand{\boxspacing}{\kern\kvtcb@left@rule\kern\kvtcb@boxsep}
    \makeatother
    \newcommand{\prompt}[4]{
        {\ttfamily\llap{{\color{#2}[#3]:\hspace{3pt}#4}}\vspace{-\baselineskip}}
    }
    

    
    % Prevent overflowing lines due to hard-to-break entities
    \sloppy
    % Setup hyperref package
    \hypersetup{
      breaklinks=true,  % so long urls are correctly broken across lines
      colorlinks=true,
      urlcolor=urlcolor,
      linkcolor=linkcolor,
      citecolor=citecolor,
      }
    % Slightly bigger margins than the latex defaults
    
    \geometry{verbose,tmargin=1in,bmargin=1in,lmargin=1in,rmargin=1in}
    
    

\usepackage{xeCJK}
\setCJKmainfont{Noto Serif CJK SC}
\setCJKmonofont{Noto Sans Mono CJK SC}
\begin{document}
    
    \maketitle
    
    

    
    \section{CS224N Assignment 1: Exploring Word Vectors (25
Points)}\label{cs224n-assignment-1-exploring-word-vectors-25-points}

\subsubsection{\texorpdfstring{ Due 4:30pm, Tue April 9th
2024}{ Due 4:30pm, Tue April 9th 2024}}\label{due-430pm-tue-april-9th-2024}

Welcome to CS224N!

Before you start, make sure you \textbf{read the README.md} in the same
directory as this notebook for important setup information. You need to
install some Python libraries before you can successfully do this
assignment. A lot of code is provided in this notebook, and we highly
encourage you to read and understand it as part of the learning :)

If you aren't super familiar with Python, Numpy, or Matplotlib, we
recommend you check out the review session on Friday. The session will
be recorded and the material will be made available on our
\href{http://web.stanford.edu/class/cs224n/index.html\#schedule}{website}.
The CS231N Python/Numpy
\href{https://cs231n.github.io/python-numpy-tutorial/}{tutorial} is also
a great resource.

\textbf{Assignment Notes:} Please make sure to save the notebook as you
go along. Submission Instructions are located at the bottom of the
notebook.

    \begin{tcolorbox}[breakable, size=fbox, boxrule=1pt, pad at break*=1mm,colback=cellbackground, colframe=cellborder]
\prompt{In}{incolor}{1}{\boxspacing}
\begin{Verbatim}[commandchars=\\\{\}]
\PY{c+c1}{\PYZsh{} All Import Statements Defined Here}
\PY{c+c1}{\PYZsh{} Note: Do not add to this list.}
\PY{c+c1}{\PYZsh{} \PYZhy{}\PYZhy{}\PYZhy{}\PYZhy{}\PYZhy{}\PYZhy{}\PYZhy{}\PYZhy{}\PYZhy{}\PYZhy{}\PYZhy{}\PYZhy{}\PYZhy{}\PYZhy{}\PYZhy{}\PYZhy{}}

\PY{k+kn}{import}\PY{+w}{ }\PY{n+nn}{sys}
\PY{k}{assert} \PY{n}{sys}\PY{o}{.}\PY{n}{version\PYZus{}info}\PY{p}{[}\PY{l+m+mi}{0}\PY{p}{]} \PY{o}{==} \PY{l+m+mi}{3}
\PY{k}{assert} \PY{n}{sys}\PY{o}{.}\PY{n}{version\PYZus{}info}\PY{p}{[}\PY{l+m+mi}{1}\PY{p}{]} \PY{o}{\PYZgt{}}\PY{o}{=} \PY{l+m+mi}{8}

\PY{k+kn}{from}\PY{+w}{ }\PY{n+nn}{platform}\PY{+w}{ }\PY{k+kn}{import} \PY{n}{python\PYZus{}version}
\PY{k}{assert} \PY{n+nb}{int}\PY{p}{(}\PY{n}{python\PYZus{}version}\PY{p}{(}\PY{p}{)}\PY{o}{.}\PY{n}{split}\PY{p}{(}\PY{l+s+s2}{\PYZdq{}}\PY{l+s+s2}{.}\PY{l+s+s2}{\PYZdq{}}\PY{p}{)}\PY{p}{[}\PY{l+m+mi}{1}\PY{p}{]}\PY{p}{)} \PY{o}{\PYZgt{}}\PY{o}{=} \PY{l+m+mi}{5}\PY{p}{,} \PY{l+s+s2}{\PYZdq{}}\PY{l+s+s2}{Please upgrade your Python version following the instructions in }\PY{l+s+se}{\PYZbs{}}
\PY{l+s+s2}{    the README.md file found in the same directory as this notebook. Your Python version is }\PY{l+s+s2}{\PYZdq{}} \PY{o}{+} \PY{n}{python\PYZus{}version}\PY{p}{(}\PY{p}{)}

\PY{k+kn}{from}\PY{+w}{ }\PY{n+nn}{gensim}\PY{n+nn}{.}\PY{n+nn}{models}\PY{+w}{ }\PY{k+kn}{import} \PY{n}{KeyedVectors}
\PY{k+kn}{from}\PY{+w}{ }\PY{n+nn}{gensim}\PY{n+nn}{.}\PY{n+nn}{test}\PY{n+nn}{.}\PY{n+nn}{utils}\PY{+w}{ }\PY{k+kn}{import} \PY{n}{datapath}
\PY{k+kn}{import}\PY{+w}{ }\PY{n+nn}{pprint}
\PY{k+kn}{import}\PY{+w}{ }\PY{n+nn}{matplotlib}\PY{n+nn}{.}\PY{n+nn}{pyplot}\PY{+w}{ }\PY{k}{as}\PY{+w}{ }\PY{n+nn}{plt}
\PY{n}{plt}\PY{o}{.}\PY{n}{rcParams}\PY{p}{[}\PY{l+s+s1}{\PYZsq{}}\PY{l+s+s1}{figure.figsize}\PY{l+s+s1}{\PYZsq{}}\PY{p}{]} \PY{o}{=} \PY{p}{[}\PY{l+m+mi}{10}\PY{p}{,} \PY{l+m+mi}{5}\PY{p}{]}

\PY{k+kn}{from}\PY{+w}{ }\PY{n+nn}{datasets}\PY{+w}{ }\PY{k+kn}{import} \PY{n}{load\PYZus{}dataset}
\PY{n}{imdb\PYZus{}dataset} \PY{o}{=} \PY{n}{load\PYZus{}dataset}\PY{p}{(}\PY{l+s+s2}{\PYZdq{}}\PY{l+s+s2}{stanfordnlp/imdb}\PY{l+s+s2}{\PYZdq{}}\PY{p}{)}

\PY{k+kn}{import}\PY{+w}{ }\PY{n+nn}{re}
\PY{k+kn}{import}\PY{+w}{ }\PY{n+nn}{numpy}\PY{+w}{ }\PY{k}{as}\PY{+w}{ }\PY{n+nn}{np}
\PY{k+kn}{import}\PY{+w}{ }\PY{n+nn}{random}
\PY{k+kn}{import}\PY{+w}{ }\PY{n+nn}{scipy}\PY{+w}{ }\PY{k}{as}\PY{+w}{ }\PY{n+nn}{sp}
\PY{k+kn}{from}\PY{+w}{ }\PY{n+nn}{sklearn}\PY{n+nn}{.}\PY{n+nn}{decomposition}\PY{+w}{ }\PY{k+kn}{import} \PY{n}{TruncatedSVD}
\PY{k+kn}{from}\PY{+w}{ }\PY{n+nn}{sklearn}\PY{n+nn}{.}\PY{n+nn}{decomposition}\PY{+w}{ }\PY{k+kn}{import} \PY{n}{PCA}

\PY{n}{START\PYZus{}TOKEN} \PY{o}{=} \PY{l+s+s1}{\PYZsq{}}\PY{l+s+s1}{\PYZlt{}START\PYZgt{}}\PY{l+s+s1}{\PYZsq{}}
\PY{n}{END\PYZus{}TOKEN} \PY{o}{=} \PY{l+s+s1}{\PYZsq{}}\PY{l+s+s1}{\PYZlt{}END\PYZgt{}}\PY{l+s+s1}{\PYZsq{}}
\PY{n}{NUM\PYZus{}SAMPLES} \PY{o}{=} \PY{l+m+mi}{150}

\PY{n}{np}\PY{o}{.}\PY{n}{random}\PY{o}{.}\PY{n}{seed}\PY{p}{(}\PY{l+m+mi}{0}\PY{p}{)}
\PY{n}{random}\PY{o}{.}\PY{n}{seed}\PY{p}{(}\PY{l+m+mi}{0}\PY{p}{)}
\PY{c+c1}{\PYZsh{} \PYZhy{}\PYZhy{}\PYZhy{}\PYZhy{}\PYZhy{}\PYZhy{}\PYZhy{}\PYZhy{}\PYZhy{}\PYZhy{}\PYZhy{}\PYZhy{}\PYZhy{}\PYZhy{}\PYZhy{}\PYZhy{}}
\end{Verbatim}
\end{tcolorbox}

    \begin{Verbatim}[commandchars=\\\{\}]
/home/zzy/study/cs224n/assignment1/a1/student/.venv/lib/python3.12/site-
packages/tqdm/auto.py:21: TqdmWarning: IProgress not found. Please update
jupyter and ipywidgets. See
https://ipywidgets.readthedocs.io/en/stable/user\_install.html
  from .autonotebook import tqdm as notebook\_tqdm
Downloading readme: 7.81kB [00:00, 7.11MB/s]
Downloading data: 100\%|██████████| 21.0M/21.0M [00:17<00:00, 1.20MB/s]
Downloading data: 100\%|██████████| 20.5M/20.5M [00:13<00:00, 1.54MB/s]
Downloading data: 100\%|██████████| 42.0M/42.0M [00:26<00:00, 1.56MB/s]
Generating train split: 100\%|██████████| 25000/25000 [00:00<00:00, 153740.45
examples/s]
Generating test split: 100\%|██████████| 25000/25000 [00:00<00:00, 155241.10
examples/s]
Generating unsupervised split: 100\%|██████████| 50000/50000 [00:00<00:00,
152516.37 examples/s]
    \end{Verbatim}

    \subsection{Word Vectors}\label{word-vectors}

Word Vectors are often used as a fundamental component for downstream
NLP tasks, e.g.~question answering, text generation, translation, etc.,
so it is important to build some intuitions as to their strengths and
weaknesses. Here, you will explore two types of word vectors: those
derived from \emph{co-occurrence matrices}, and those derived via
\emph{GloVe}.

\textbf{Note on Terminology:} The terms ``word vectors'' and ``word
embeddings'' are often used interchangeably. The term ``embedding''
refers to the fact that we are encoding aspects of a word's meaning in a
lower dimensional space. As
\href{https://en.wikipedia.org/wiki/Word_embedding}{Wikipedia} states,
``\emph{conceptually it involves a mathematical embedding from a space
with one dimension per word to a continuous vector space with a much
lower dimension}''.

    \subsection{Part 1: Count-Based Word Vectors (10
points)}\label{part-1-count-based-word-vectors-10-points}

Most word vector models start from the following idea:

\emph{You shall know a word by the company it keeps
(\href{https://en.wikipedia.org/wiki/John_Rupert_Firth}{Firth, J. R.
1957:11})}

Many word vector implementations are driven by the idea that similar
words, i.e., (near) synonyms, will be used in similar contexts. As a
result, similar words will often be spoken or written along with a
shared subset of words, i.e., contexts. By examining these contexts, we
can try to develop embeddings for our words. With this intuition in
mind, many ``old school'' approaches to constructing word vectors relied
on word counts. Here we elaborate upon one of those strategies,
\emph{co-occurrence matrices} (for more information, see
\href{https://web.stanford.edu/~jurafsky/slp3/6.pdf}{here} or
\href{https://web.archive.org/web/20190530091127/https://medium.com/data-science-group-iitr/word-embedding-2d05d270b285}{here}).

    \subsubsection{Co-Occurrence}\label{co-occurrence}

A co-occurrence matrix counts how often things co-occur in some
environment. Given some word \(w_i\) occurring in the document, we
consider the \emph{context window} surrounding \(w_i\). Supposing our
fixed window size is \(n\), then this is the \(n\) preceding and \(n\)
subsequent words in that document, i.e.~words \(w_{i-n} \dots w_{i-1}\)
and \(w_{i+1} \dots w_{i+n}\). We build a \emph{co-occurrence matrix}
\(M\), which is a symmetric word-by-word matrix in which \(M_{ij}\) is
the number of times \(w_j\) appears inside \(w_i\)'s window among all
documents.

\textbf{Example: Co-Occurrence with Fixed Window of n=1}:

Document 1: ``all that glitters is not gold''

Document 2: ``all is well that ends well''

\begin{longtable}[]{@{}
  >{\raggedright\arraybackslash}p{(\linewidth - 20\tabcolsep) * \real{0.1351}}
  >{\raggedright\arraybackslash}p{(\linewidth - 20\tabcolsep) * \real{0.0946}}
  >{\raggedright\arraybackslash}p{(\linewidth - 20\tabcolsep) * \real{0.0676}}
  >{\raggedright\arraybackslash}p{(\linewidth - 20\tabcolsep) * \real{0.0811}}
  >{\raggedright\arraybackslash}p{(\linewidth - 20\tabcolsep) * \real{0.1351}}
  >{\raggedright\arraybackslash}p{(\linewidth - 20\tabcolsep) * \real{0.0811}}
  >{\raggedright\arraybackslash}p{(\linewidth - 20\tabcolsep) * \real{0.0811}}
  >{\raggedright\arraybackslash}p{(\linewidth - 20\tabcolsep) * \real{0.0946}}
  >{\raggedright\arraybackslash}p{(\linewidth - 20\tabcolsep) * \real{0.0811}}
  >{\raggedright\arraybackslash}p{(\linewidth - 20\tabcolsep) * \real{0.0811}}
  >{\raggedright\arraybackslash}p{(\linewidth - 20\tabcolsep) * \real{0.0676}}@{}}
\toprule\noalign{}
\begin{minipage}[b]{\linewidth}\raggedright
*
\end{minipage} & \begin{minipage}[b]{\linewidth}\raggedright
\texttt{\textless{}START\textgreater{}}
\end{minipage} & \begin{minipage}[b]{\linewidth}\raggedright
all
\end{minipage} & \begin{minipage}[b]{\linewidth}\raggedright
that
\end{minipage} & \begin{minipage}[b]{\linewidth}\raggedright
glitters
\end{minipage} & \begin{minipage}[b]{\linewidth}\raggedright
is
\end{minipage} & \begin{minipage}[b]{\linewidth}\raggedright
not
\end{minipage} & \begin{minipage}[b]{\linewidth}\raggedright
gold
\end{minipage} & \begin{minipage}[b]{\linewidth}\raggedright
well
\end{minipage} & \begin{minipage}[b]{\linewidth}\raggedright
ends
\end{minipage} & \begin{minipage}[b]{\linewidth}\raggedright
\texttt{\textless{}END\textgreater{}}
\end{minipage} \\
\midrule\noalign{}
\endhead
\bottomrule\noalign{}
\endlastfoot
\texttt{\textless{}START\textgreater{}} & 0 & 2 & 0 & 0 & 0 & 0 & 0 & 0
& 0 & 0 \\
all & 2 & 0 & 1 & 0 & 1 & 0 & 0 & 0 & 0 & 0 \\
that & 0 & 1 & 0 & 1 & 0 & 0 & 0 & 1 & 1 & 0 \\
glitters & 0 & 0 & 1 & 0 & 1 & 0 & 0 & 0 & 0 & 0 \\
is & 0 & 1 & 0 & 1 & 0 & 1 & 0 & 1 & 0 & 0 \\
not & 0 & 0 & 0 & 0 & 1 & 0 & 1 & 0 & 0 & 0 \\
gold & 0 & 0 & 0 & 0 & 0 & 1 & 0 & 0 & 0 & 1 \\
well & 0 & 0 & 1 & 0 & 1 & 0 & 0 & 0 & 1 & 1 \\
ends & 0 & 0 & 1 & 0 & 0 & 0 & 0 & 1 & 0 & 0 \\
\texttt{\textless{}END\textgreater{}} & 0 & 0 & 0 & 0 & 0 & 0 & 1 & 1 &
0 & 0 \\
\end{longtable}

In NLP, we commonly use \texttt{\textless{}START\textgreater{}} and
\texttt{\textless{}END\textgreater{}} tokens to mark the beginning and
end of sentences, paragraphs, or documents. These tokens are included in
co-occurrence counts, encapsulating each document, for example:
``\texttt{\textless{}START\textgreater{}} All that glitters is not gold
\texttt{\textless{}END\textgreater{}}''.

The matrix rows (or columns) provide word vectors based on word-word
co-occurrence, but they can be large. To reduce dimensionality, we
employ Singular Value Decomposition (SVD), akin to PCA, selecting the
top \(k\) principal components. The SVD process decomposes the
co-occurrence matrix \(A\) into singular values in the diagonal \(S\)
matrix and new, shorter word vectors in \(U_k\).

This dimensionality reduction maintains semantic relationships; for
instance, \emph{doctor} and \emph{hospital} will be closer than
\emph{doctor} and \emph{dog}.

For those unfamiliar with eigenvalues and SVD, a beginner-friendly
introduction to SVD is available
\href{https://davetang.org/file/Singular_Value_Decomposition_Tutorial.pdf}{here}.
Additional resources for in-depth understanding include lectures
\href{https://web.stanford.edu/class/cs168/l/l7.pdf}{7},
\href{http://theory.stanford.edu/~tim/s15/l/l8.pdf}{8}, and
\href{https://web.stanford.edu/class/cs168/l/l9.pdf}{9} of CS168,
providing high-level treatment of these algorithms. For practical
implementation, utilizing pre-programmed functions from Python packages
like numpy, scipy, or sklearn is recommended. While applying full SVD to
large corpora can be memory-intensive, scalable techniques such as
Truncated SVD exist for extracting the top \(k\) vector components
efficiently.

    \subsubsection{Plotting Co-Occurrence Word
Embeddings}\label{plotting-co-occurrence-word-embeddings}

Here, we will be using the Large Movie Review Dataset. This is a dataset
for binary sentiment classification containing substantially more data
than previous benchmark datasets. We provide a set of 25,000 highly
polar movie reviews for training, and 25,000 for testing. There is
additional unlabeled data for use as well. We provide a
\texttt{read\_corpus} function below that pulls out the text of a movie
review from the dataset. The function also adds
\texttt{\textless{}START\textgreater{}} and
\texttt{\textless{}END\textgreater{}} tokens to each of the documents,
and lowercases words. You do \textbf{not} have to perform any other kind
of pre-processing.

    \begin{tcolorbox}[breakable, size=fbox, boxrule=1pt, pad at break*=1mm,colback=cellbackground, colframe=cellborder]
\prompt{In}{incolor}{2}{\boxspacing}
\begin{Verbatim}[commandchars=\\\{\}]
\PY{k}{def}\PY{+w}{ }\PY{n+nf}{read\PYZus{}corpus}\PY{p}{(}\PY{p}{)}\PY{p}{:}
\PY{+w}{    }\PY{l+s+sd}{\PYZdq{}\PYZdq{}\PYZdq{} Read files from the Large Movie Review Dataset.}
\PY{l+s+sd}{        Params:}
\PY{l+s+sd}{            category (string): category name}
\PY{l+s+sd}{        Return:}
\PY{l+s+sd}{            list of lists, with words from each of the processed files}
\PY{l+s+sd}{    \PYZdq{}\PYZdq{}\PYZdq{}}
    \PY{n}{files} \PY{o}{=} \PY{n}{imdb\PYZus{}dataset}\PY{p}{[}\PY{l+s+s2}{\PYZdq{}}\PY{l+s+s2}{train}\PY{l+s+s2}{\PYZdq{}}\PY{p}{]}\PY{p}{[}\PY{l+s+s2}{\PYZdq{}}\PY{l+s+s2}{text}\PY{l+s+s2}{\PYZdq{}}\PY{p}{]}\PY{p}{[}\PY{p}{:}\PY{n}{NUM\PYZus{}SAMPLES}\PY{p}{]}
    \PY{k}{return} \PY{p}{[}\PY{p}{[}\PY{n}{START\PYZus{}TOKEN}\PY{p}{]} \PY{o}{+} \PY{p}{[}\PY{n}{re}\PY{o}{.}\PY{n}{sub}\PY{p}{(}\PY{l+s+sa}{r}\PY{l+s+s1}{\PYZsq{}}\PY{l+s+s1}{[\PYZca{}}\PY{l+s+s1}{\PYZbs{}}\PY{l+s+s1}{w]}\PY{l+s+s1}{\PYZsq{}}\PY{p}{,} \PY{l+s+s1}{\PYZsq{}}\PY{l+s+s1}{\PYZsq{}}\PY{p}{,} \PY{n}{w}\PY{o}{.}\PY{n}{lower}\PY{p}{(}\PY{p}{)}\PY{p}{)} \PY{k}{for} \PY{n}{w} \PY{o+ow}{in} \PY{n}{f}\PY{o}{.}\PY{n}{split}\PY{p}{(}\PY{l+s+s2}{\PYZdq{}}\PY{l+s+s2}{ }\PY{l+s+s2}{\PYZdq{}}\PY{p}{)}\PY{p}{]} \PY{o}{+} \PY{p}{[}\PY{n}{END\PYZus{}TOKEN}\PY{p}{]} \PY{k}{for} \PY{n}{f} \PY{o+ow}{in} \PY{n}{files}\PY{p}{]}
\end{Verbatim}
\end{tcolorbox}

    Let's have a look what these documents are like\ldots.

    \begin{tcolorbox}[breakable, size=fbox, boxrule=1pt, pad at break*=1mm,colback=cellbackground, colframe=cellborder]
\prompt{In}{incolor}{3}{\boxspacing}
\begin{Verbatim}[commandchars=\\\{\}]
\PY{n}{imdb\PYZus{}corpus} \PY{o}{=} \PY{n}{read\PYZus{}corpus}\PY{p}{(}\PY{p}{)}
\PY{n}{pprint}\PY{o}{.}\PY{n}{pprint}\PY{p}{(}\PY{n}{imdb\PYZus{}corpus}\PY{p}{[}\PY{p}{:}\PY{l+m+mi}{3}\PY{p}{]}\PY{p}{,} \PY{n}{compact}\PY{o}{=}\PY{k+kc}{True}\PY{p}{,} \PY{n}{width}\PY{o}{=}\PY{l+m+mi}{100}\PY{p}{)}
\PY{n+nb}{print}\PY{p}{(}\PY{l+s+s2}{\PYZdq{}}\PY{l+s+s2}{corpus size: }\PY{l+s+s2}{\PYZdq{}}\PY{p}{,} \PY{n+nb}{len}\PY{p}{(}\PY{n}{imdb\PYZus{}corpus}\PY{p}{[}\PY{l+m+mi}{0}\PY{p}{]}\PY{p}{)}\PY{p}{)}
\end{Verbatim}
\end{tcolorbox}

    \begin{Verbatim}[commandchars=\\\{\}]
[['<START>', 'i', 'rented', 'i', 'am', 'curiousyellow', 'from', 'my', 'video',
'store', 'because',
  'of', 'all', 'the', 'controversy', 'that', 'surrounded', 'it', 'when', 'it',
'was', 'first',
  'released', 'in', '1967', 'i', 'also', 'heard', 'that', 'at', 'first', 'it',
'was', 'seized',
  'by', 'us', 'customs', 'if', 'it', 'ever', 'tried', 'to', 'enter', 'this',
'country', 'therefore',
  'being', 'a', 'fan', 'of', 'films', 'considered', 'controversial', 'i',
'really', 'had', 'to',
  'see', 'this', 'for', 'myselfbr', 'br', 'the', 'plot', 'is', 'centered',
'around', 'a', 'young',
  'swedish', 'drama', 'student', 'named', 'lena', 'who', 'wants', 'to', 'learn',
'everything',
  'she', 'can', 'about', 'life', 'in', 'particular', 'she', 'wants', 'to',
'focus', 'her',
  'attentions', 'to', 'making', 'some', 'sort', 'of', 'documentary', 'on',
'what', 'the', 'average',
  'swede', 'thought', 'about', 'certain', 'political', 'issues', 'such', 'as',
'the', 'vietnam',
  'war', 'and', 'race', 'issues', 'in', 'the', 'united', 'states', 'in',
'between', 'asking',
  'politicians', 'and', 'ordinary', 'denizens', 'of', 'stockholm', 'about',
'their', 'opinions',
  'on', 'politics', 'she', 'has', 'sex', 'with', 'her', 'drama', 'teacher',
'classmates', 'and',
  'married', 'menbr', 'br', 'what', 'kills', 'me', 'about', 'i', 'am',
'curiousyellow', 'is',
  'that', '40', 'years', 'ago', 'this', 'was', 'considered', 'pornographic',
'really', 'the', 'sex',
  'and', 'nudity', 'scenes', 'are', 'few', 'and', 'far', 'between', 'even',
'then', 'its', 'not',
  'shot', 'like', 'some', 'cheaply', 'made', 'porno', 'while', 'my',
'countrymen', 'mind', 'find',
  'it', 'shocking', 'in', 'reality', 'sex', 'and', 'nudity', 'are', 'a',
'major', 'staple', 'in',
  'swedish', 'cinema', 'even', 'ingmar', 'bergman', 'arguably', 'their',
'answer', 'to', 'good',
  'old', 'boy', 'john', 'ford', 'had', 'sex', 'scenes', 'in', 'his', 'filmsbr',
'br', 'i', 'do',
  'commend', 'the', 'filmmakers', 'for', 'the', 'fact', 'that', 'any', 'sex',
'shown', 'in', 'the',
  'film', 'is', 'shown', 'for', 'artistic', 'purposes', 'rather', 'than',
'just', 'to', 'shock',
  'people', 'and', 'make', 'money', 'to', 'be', 'shown', 'in', 'pornographic',
'theaters', 'in',
  'america', 'i', 'am', 'curiousyellow', 'is', 'a', 'good', 'film', 'for',
'anyone', 'wanting',
  'to', 'study', 'the', 'meat', 'and', 'potatoes', 'no', 'pun', 'intended',
'of', 'swedish',
  'cinema', 'but', 'really', 'this', 'film', 'doesnt', 'have', 'much', 'of',
'a', 'plot', '<END>'],
 ['<START>', 'i', 'am', 'curious', 'yellow', 'is', 'a', 'risible', 'and',
'pretentious', 'steaming',
  'pile', 'it', 'doesnt', 'matter', 'what', 'ones', 'political', 'views', 'are',
'because', 'this',
  'film', 'can', 'hardly', 'be', 'taken', 'seriously', 'on', 'any', 'level',
'as', 'for', 'the',
  'claim', 'that', 'frontal', 'male', 'nudity', 'is', 'an', 'automatic', 'nc17',
'that', 'isnt',
  'true', 'ive', 'seen', 'rrated', 'films', 'with', 'male', 'nudity', 'granted',
'they', 'only',
  'offer', 'some', 'fleeting', 'views', 'but', 'where', 'are', 'the', 'rrated',
'films', 'with',
  'gaping', 'vulvas', 'and', 'flapping', 'labia', 'nowhere', 'because', 'they',
'dont', 'exist',
  'the', 'same', 'goes', 'for', 'those', 'crappy', 'cable', 'shows', 'schlongs',
'swinging', 'in',
  'the', 'breeze', 'but', 'not', 'a', 'clitoris', 'in', 'sight', 'and', 'those',
'pretentious',
  'indie', 'movies', 'like', 'the', 'brown', 'bunny', 'in', 'which', 'were',
'treated', 'to', 'the',
  'site', 'of', 'vincent', 'gallos', 'throbbing', 'johnson', 'but', 'not', 'a',
'trace', 'of',
  'pink', 'visible', 'on', 'chloe', 'sevigny', 'before', 'crying', 'or',
'implying',
  'doublestandard', 'in', 'matters', 'of', 'nudity', 'the', 'mentally',
'obtuse', 'should', 'take',
  'into', 'account', 'one', 'unavoidably', 'obvious', 'anatomical',
'difference', 'between', 'men',
  'and', 'women', 'there', 'are', 'no', 'genitals', 'on', 'display', 'when',
'actresses', 'appears',
  'nude', 'and', 'the', 'same', 'cannot', 'be', 'said', 'for', 'a', 'man', 'in',
'fact', 'you',
  'generally', 'wont', 'see', 'female', 'genitals', 'in', 'an', 'american',
'film', 'in',
  'anything', 'short', 'of', 'porn', 'or', 'explicit', 'erotica', 'this',
'alleged',
  'doublestandard', 'is', 'less', 'a', 'double', 'standard', 'than', 'an',
'admittedly',
  'depressing', 'ability', 'to', 'come', 'to', 'terms', 'culturally', 'with',
'the', 'insides',
  'of', 'womens', 'bodies', '<END>'],
 ['<START>', 'if', 'only', 'to', 'avoid', 'making', 'this', 'type', 'of',
'film', 'in', 'the',
  'future', 'this', 'film', 'is', 'interesting', 'as', 'an', 'experiment',
'but', 'tells', 'no',
  'cogent', 'storybr', 'br', 'one', 'might', 'feel', 'virtuous', 'for',
'sitting', 'thru', 'it',
  'because', 'it', 'touches', 'on', 'so', 'many', 'important', 'issues', 'but',
'it', 'does', 'so',
  'without', 'any', 'discernable', 'motive', 'the', 'viewer', 'comes', 'away',
'with', 'no', 'new',
  'perspectives', 'unless', 'one', 'comes', 'up', 'with', 'one', 'while',
'ones', 'mind', 'wanders',
  'as', 'it', 'will', 'invariably', 'do', 'during', 'this', 'pointless',
'filmbr', 'br', 'one',
  'might', 'better', 'spend', 'ones', 'time', 'staring', 'out', 'a', 'window',
'at', 'a', 'tree',
  'growingbr', 'br', '', '<END>']]
corpus size:  290
    \end{Verbatim}

    \subsubsection{\texorpdfstring{Question 1.1: Implement
\texttt{distinct\_words} {[}code{]} (2
points)}{Question 1.1: Implement distinct\_words {[}code{]} (2 points)}}\label{question-1.1-implement-distinct_words-code-2-points}

Write a method to work out the distinct words (word types) that occur in
the corpus.

You can use \texttt{for} loops to process the input \texttt{corpus} (a
list of list of strings), but try using Python list comprehensions
(which are generally faster). In particular,
\href{https://coderwall.com/p/rcmaea/flatten-a-list-of-lists-in-one-line-in-python}{this}
may be useful to flatten a list of lists. If you're not familiar with
Python list comprehensions in general, here's
\href{https://python-3-patterns-idioms-test.readthedocs.io/en/latest/Comprehensions.html}{more
information}.

Your returned \texttt{corpus\_words} should be sorted. You can use
python's \texttt{sorted} function for this.

You may find it useful to use
\href{https://www.w3schools.com/python/python_sets.asp}{Python sets} to
remove duplicate words.

    \begin{tcolorbox}[breakable, size=fbox, boxrule=1pt, pad at break*=1mm,colback=cellbackground, colframe=cellborder]
\prompt{In}{incolor}{ }{\boxspacing}
\begin{Verbatim}[commandchars=\\\{\}]
\PY{k}{def}\PY{+w}{ }\PY{n+nf}{distinct\PYZus{}words}\PY{p}{(}\PY{n}{corpus}\PY{p}{)}\PY{p}{:}
\PY{+w}{    }\PY{l+s+sd}{\PYZdq{}\PYZdq{}\PYZdq{} Determine a list of distinct words for the corpus.}
\PY{l+s+sd}{        Params:}
\PY{l+s+sd}{            corpus (list of list of strings): corpus of documents}
\PY{l+s+sd}{        Return:}
\PY{l+s+sd}{            corpus\PYZus{}words (list of strings): sorted list of distinct words across the corpus}
\PY{l+s+sd}{            n\PYZus{}corpus\PYZus{}words (integer): number of distinct words across the corpus}
\PY{l+s+sd}{    \PYZdq{}\PYZdq{}\PYZdq{}}
    \PY{n}{corpus\PYZus{}words} \PY{o}{=} \PY{p}{[}\PY{p}{]}
    \PY{n}{n\PYZus{}corpus\PYZus{}words} \PY{o}{=} \PY{o}{\PYZhy{}}\PY{l+m+mi}{1}

    \PY{c+c1}{\PYZsh{} \PYZhy{}\PYZhy{}\PYZhy{}\PYZhy{}\PYZhy{}\PYZhy{}\PYZhy{}\PYZhy{}\PYZhy{}\PYZhy{}\PYZhy{}\PYZhy{}\PYZhy{}\PYZhy{}\PYZhy{}\PYZhy{}\PYZhy{}\PYZhy{}}
    \PY{c+c1}{\PYZsh{} Write your implementation here.}
    \PY{k}{for} \PY{n}{doc} \PY{o+ow}{in} \PY{n}{corpus}\PY{p}{:}
        \PY{k}{for} \PY{n}{word} \PY{o+ow}{in} \PY{n}{doc}\PY{p}{:}
            \PY{k}{if} \PY{n}{word} \PY{o+ow}{not} \PY{o+ow}{in} \PY{n}{corpus\PYZus{}words}\PY{p}{:}
                \PY{n}{corpus\PYZus{}words}\PY{o}{.}\PY{n}{append}\PY{p}{(}\PY{n}{word}\PY{p}{)}

    \PY{c+c1}{\PYZsh{} \PYZhy{}\PYZhy{}\PYZhy{}\PYZhy{}\PYZhy{}\PYZhy{}\PYZhy{}\PYZhy{}\PYZhy{}\PYZhy{}\PYZhy{}\PYZhy{}\PYZhy{}\PYZhy{}\PYZhy{}\PYZhy{}\PYZhy{}\PYZhy{}}
    \PY{n}{corpus\PYZus{}words}\PY{o}{.}\PY{n}{sort}\PY{p}{(}\PY{p}{)}
    \PY{n}{n\PYZus{}corpus\PYZus{}words} \PY{o}{=} \PY{n+nb}{len}\PY{p}{(}\PY{n}{corpus\PYZus{}words}\PY{p}{)}
    \PY{c+c1}{\PYZsh{} \PYZhy{}\PYZhy{}\PYZhy{}\PYZhy{}\PYZhy{}\PYZhy{}\PYZhy{}\PYZhy{}\PYZhy{}\PYZhy{}\PYZhy{}\PYZhy{}\PYZhy{}\PYZhy{}\PYZhy{}\PYZhy{}\PYZhy{}\PYZhy{}}

    \PY{k}{return} \PY{n}{corpus\PYZus{}words}\PY{p}{,} \PY{n}{n\PYZus{}corpus\PYZus{}words}
\end{Verbatim}
\end{tcolorbox}

    \begin{tcolorbox}[breakable, size=fbox, boxrule=1pt, pad at break*=1mm,colback=cellbackground, colframe=cellborder]
\prompt{In}{incolor}{9}{\boxspacing}
\begin{Verbatim}[commandchars=\\\{\}]
\PY{c+c1}{\PYZsh{} \PYZhy{}\PYZhy{}\PYZhy{}\PYZhy{}\PYZhy{}\PYZhy{}\PYZhy{}\PYZhy{}\PYZhy{}\PYZhy{}\PYZhy{}\PYZhy{}\PYZhy{}\PYZhy{}\PYZhy{}\PYZhy{}\PYZhy{}\PYZhy{}\PYZhy{}\PYZhy{}\PYZhy{}}
\PY{c+c1}{\PYZsh{} Run this sanity check}
\PY{c+c1}{\PYZsh{} Note that this not an exhaustive check for correctness.}
\PY{c+c1}{\PYZsh{} \PYZhy{}\PYZhy{}\PYZhy{}\PYZhy{}\PYZhy{}\PYZhy{}\PYZhy{}\PYZhy{}\PYZhy{}\PYZhy{}\PYZhy{}\PYZhy{}\PYZhy{}\PYZhy{}\PYZhy{}\PYZhy{}\PYZhy{}\PYZhy{}\PYZhy{}\PYZhy{}\PYZhy{}}

\PY{c+c1}{\PYZsh{} Define toy corpus}
\PY{n}{test\PYZus{}corpus} \PY{o}{=} \PY{p}{[}\PY{l+s+s2}{\PYZdq{}}\PY{l+s+si}{\PYZob{}\PYZcb{}}\PY{l+s+s2}{ All that glitters isn}\PY{l+s+s2}{\PYZsq{}}\PY{l+s+s2}{t gold }\PY{l+s+si}{\PYZob{}\PYZcb{}}\PY{l+s+s2}{\PYZdq{}}\PY{o}{.}\PY{n}{format}\PY{p}{(}\PY{n}{START\PYZus{}TOKEN}\PY{p}{,} \PY{n}{END\PYZus{}TOKEN}\PY{p}{)}\PY{o}{.}\PY{n}{split}\PY{p}{(}\PY{l+s+s2}{\PYZdq{}}\PY{l+s+s2}{ }\PY{l+s+s2}{\PYZdq{}}\PY{p}{)}\PY{p}{,} \PY{l+s+s2}{\PYZdq{}}\PY{l+s+si}{\PYZob{}\PYZcb{}}\PY{l+s+s2}{ All}\PY{l+s+s2}{\PYZsq{}}\PY{l+s+s2}{s well that ends well }\PY{l+s+si}{\PYZob{}\PYZcb{}}\PY{l+s+s2}{\PYZdq{}}\PY{o}{.}\PY{n}{format}\PY{p}{(}\PY{n}{START\PYZus{}TOKEN}\PY{p}{,} \PY{n}{END\PYZus{}TOKEN}\PY{p}{)}\PY{o}{.}\PY{n}{split}\PY{p}{(}\PY{l+s+s2}{\PYZdq{}}\PY{l+s+s2}{ }\PY{l+s+s2}{\PYZdq{}}\PY{p}{)}\PY{p}{]}
\PY{n}{test\PYZus{}corpus\PYZus{}words}\PY{p}{,} \PY{n}{num\PYZus{}corpus\PYZus{}words} \PY{o}{=} \PY{n}{distinct\PYZus{}words}\PY{p}{(}\PY{n}{test\PYZus{}corpus}\PY{p}{)}

\PY{c+c1}{\PYZsh{} Correct answers}
\PY{n}{ans\PYZus{}test\PYZus{}corpus\PYZus{}words} \PY{o}{=} \PY{n+nb}{sorted}\PY{p}{(}\PY{p}{[}\PY{n}{START\PYZus{}TOKEN}\PY{p}{,} \PY{l+s+s2}{\PYZdq{}}\PY{l+s+s2}{All}\PY{l+s+s2}{\PYZdq{}}\PY{p}{,} \PY{l+s+s2}{\PYZdq{}}\PY{l+s+s2}{ends}\PY{l+s+s2}{\PYZdq{}}\PY{p}{,} \PY{l+s+s2}{\PYZdq{}}\PY{l+s+s2}{that}\PY{l+s+s2}{\PYZdq{}}\PY{p}{,} \PY{l+s+s2}{\PYZdq{}}\PY{l+s+s2}{gold}\PY{l+s+s2}{\PYZdq{}}\PY{p}{,} \PY{l+s+s2}{\PYZdq{}}\PY{l+s+s2}{All}\PY{l+s+s2}{\PYZsq{}}\PY{l+s+s2}{s}\PY{l+s+s2}{\PYZdq{}}\PY{p}{,} \PY{l+s+s2}{\PYZdq{}}\PY{l+s+s2}{glitters}\PY{l+s+s2}{\PYZdq{}}\PY{p}{,} \PY{l+s+s2}{\PYZdq{}}\PY{l+s+s2}{isn}\PY{l+s+s2}{\PYZsq{}}\PY{l+s+s2}{t}\PY{l+s+s2}{\PYZdq{}}\PY{p}{,} \PY{l+s+s2}{\PYZdq{}}\PY{l+s+s2}{well}\PY{l+s+s2}{\PYZdq{}}\PY{p}{,} \PY{n}{END\PYZus{}TOKEN}\PY{p}{]}\PY{p}{)}
\PY{n}{ans\PYZus{}num\PYZus{}corpus\PYZus{}words} \PY{o}{=} \PY{n+nb}{len}\PY{p}{(}\PY{n}{ans\PYZus{}test\PYZus{}corpus\PYZus{}words}\PY{p}{)}

\PY{c+c1}{\PYZsh{} Test correct number of words}
\PY{k}{assert}\PY{p}{(}\PY{n}{num\PYZus{}corpus\PYZus{}words} \PY{o}{==} \PY{n}{ans\PYZus{}num\PYZus{}corpus\PYZus{}words}\PY{p}{)}\PY{p}{,} \PY{l+s+s2}{\PYZdq{}}\PY{l+s+s2}{Incorrect number of distinct words. Correct: }\PY{l+s+si}{\PYZob{}\PYZcb{}}\PY{l+s+s2}{. Yours: }\PY{l+s+si}{\PYZob{}\PYZcb{}}\PY{l+s+s2}{\PYZdq{}}\PY{o}{.}\PY{n}{format}\PY{p}{(}\PY{n}{ans\PYZus{}num\PYZus{}corpus\PYZus{}words}\PY{p}{,} \PY{n}{num\PYZus{}corpus\PYZus{}words}\PY{p}{)}

\PY{c+c1}{\PYZsh{} Test correct words}
\PY{k}{assert} \PY{p}{(}\PY{n}{test\PYZus{}corpus\PYZus{}words} \PY{o}{==} \PY{n}{ans\PYZus{}test\PYZus{}corpus\PYZus{}words}\PY{p}{)}\PY{p}{,} \PY{l+s+s2}{\PYZdq{}}\PY{l+s+s2}{Incorrect corpus\PYZus{}words.}\PY{l+s+se}{\PYZbs{}n}\PY{l+s+s2}{Correct: }\PY{l+s+si}{\PYZob{}\PYZcb{}}\PY{l+s+se}{\PYZbs{}n}\PY{l+s+s2}{Yours:   }\PY{l+s+si}{\PYZob{}\PYZcb{}}\PY{l+s+s2}{\PYZdq{}}\PY{o}{.}\PY{n}{format}\PY{p}{(}\PY{n+nb}{str}\PY{p}{(}\PY{n}{ans\PYZus{}test\PYZus{}corpus\PYZus{}words}\PY{p}{)}\PY{p}{,} \PY{n+nb}{str}\PY{p}{(}\PY{n}{test\PYZus{}corpus\PYZus{}words}\PY{p}{)}\PY{p}{)}

\PY{c+c1}{\PYZsh{} Print Success}
\PY{n+nb}{print} \PY{p}{(}\PY{l+s+s2}{\PYZdq{}}\PY{l+s+s2}{\PYZhy{}}\PY{l+s+s2}{\PYZdq{}} \PY{o}{*} \PY{l+m+mi}{80}\PY{p}{)}
\PY{n+nb}{print}\PY{p}{(}\PY{l+s+s2}{\PYZdq{}}\PY{l+s+s2}{Passed All Tests!}\PY{l+s+s2}{\PYZdq{}}\PY{p}{)}
\PY{n+nb}{print} \PY{p}{(}\PY{l+s+s2}{\PYZdq{}}\PY{l+s+s2}{\PYZhy{}}\PY{l+s+s2}{\PYZdq{}} \PY{o}{*} \PY{l+m+mi}{80}\PY{p}{)}
\end{Verbatim}
\end{tcolorbox}

    \begin{Verbatim}[commandchars=\\\{\}]
--------------------------------------------------------------------------------
Passed All Tests!
--------------------------------------------------------------------------------
    \end{Verbatim}

    \subsubsection{\texorpdfstring{Question 1.2: Implement
\texttt{compute\_co\_occurrence\_matrix} {[}code{]} (3
points)}{Question 1.2: Implement compute\_co\_occurrence\_matrix {[}code{]} (3 points)}}\label{question-1.2-implement-compute_co_occurrence_matrix-code-3-points}

Write a method that constructs a co-occurrence matrix for a certain
window-size \(n\) (with a default of 4), considering words \(n\) before
and \(n\) after the word in the center of the window. Here, we start to
use \texttt{numpy\ (np)} to represent vectors, matrices, and tensors. If
you're not familiar with NumPy, there's a NumPy tutorial in the second
half of this cs231n
\href{http://cs231n.github.io/python-numpy-tutorial/}{Python NumPy
tutorial}.

    \begin{tcolorbox}[breakable, size=fbox, boxrule=1pt, pad at break*=1mm,colback=cellbackground, colframe=cellborder]
\prompt{In}{incolor}{ }{\boxspacing}
\begin{Verbatim}[commandchars=\\\{\}]
\PY{k}{def}\PY{+w}{ }\PY{n+nf}{compute\PYZus{}co\PYZus{}occurrence\PYZus{}matrix}\PY{p}{(}\PY{n}{corpus}\PY{p}{,} \PY{n}{window\PYZus{}size}\PY{o}{=}\PY{l+m+mi}{4}\PY{p}{)}\PY{p}{:}
\PY{+w}{    }\PY{l+s+sd}{\PYZdq{}\PYZdq{}\PYZdq{} Compute co\PYZhy{}occurrence matrix for the given corpus and window\PYZus{}size (default of 4).}

\PY{l+s+sd}{        Note: Each word in a document should be at the center of a window. Words near edges will have a smaller}
\PY{l+s+sd}{              number of co\PYZhy{}occurring words.}

\PY{l+s+sd}{              For example, if we take the document \PYZdq{}\PYZlt{}START\PYZgt{} All that glitters is not gold \PYZlt{}END\PYZgt{}\PYZdq{} with window size of 4,}
\PY{l+s+sd}{              \PYZdq{}All\PYZdq{} will co\PYZhy{}occur with \PYZdq{}\PYZlt{}START\PYZgt{}\PYZdq{}, \PYZdq{}that\PYZdq{}, \PYZdq{}glitters\PYZdq{}, \PYZdq{}is\PYZdq{}, and \PYZdq{}not\PYZdq{}.}

\PY{l+s+sd}{        Params:}
\PY{l+s+sd}{            corpus (list of list of strings): corpus of documents}
\PY{l+s+sd}{            window\PYZus{}size (int): size of context window}
\PY{l+s+sd}{        Return:}
\PY{l+s+sd}{            M (a symmetric numpy matrix of shape (number of unique words in the corpus , number of unique words in the corpus)):}
\PY{l+s+sd}{                Co\PYZhy{}occurence matrix of word counts.}
\PY{l+s+sd}{                The ordering of the words in the rows/columns should be the same as the ordering of the words given by the distinct\PYZus{}words function.}
\PY{l+s+sd}{            word2ind (dict): dictionary that maps word to index (i.e. row/column number) for matrix M.}
\PY{l+s+sd}{    \PYZdq{}\PYZdq{}\PYZdq{}}
    \PY{n}{words}\PY{p}{,} \PY{n}{n\PYZus{}words} \PY{o}{=} \PY{n}{distinct\PYZus{}words}\PY{p}{(}\PY{n}{corpus}\PY{p}{)}
    \PY{n}{M} \PY{o}{=} \PY{k+kc}{None}
    \PY{n}{word2ind} \PY{o}{=} \PY{p}{\PYZob{}}\PY{p}{\PYZcb{}}

    \PY{c+c1}{\PYZsh{} \PYZhy{}\PYZhy{}\PYZhy{}\PYZhy{}\PYZhy{}\PYZhy{}\PYZhy{}\PYZhy{}\PYZhy{}\PYZhy{}\PYZhy{}\PYZhy{}\PYZhy{}\PYZhy{}\PYZhy{}\PYZhy{}\PYZhy{}\PYZhy{}}
    \PY{c+c1}{\PYZsh{} Write your implementation here.}
    \PY{n}{M} \PY{o}{=} \PY{n}{np}\PY{o}{.}\PY{n}{zeros}\PY{p}{(}\PY{p}{(}\PY{n}{n\PYZus{}words}\PY{p}{,} \PY{n}{n\PYZus{}words}\PY{p}{)}\PY{p}{)}
    \PY{k}{for} \PY{n}{i}\PY{p}{,} \PY{n}{word} \PY{o+ow}{in} \PY{n+nb}{enumerate}\PY{p}{(}\PY{n}{words}\PY{p}{)}\PY{p}{:}
        \PY{n}{word2ind}\PY{p}{[}\PY{n}{word}\PY{p}{]} \PY{o}{=} \PY{n}{i}
    \PY{k}{for} \PY{n}{doc} \PY{o+ow}{in} \PY{n}{corpus}\PY{p}{:}
        \PY{k}{for} \PY{n}{i}\PY{p}{,}\PY{n}{word} \PY{o+ow}{in} \PY{n+nb}{enumerate}\PY{p}{(}\PY{n}{doc}\PY{p}{)}\PY{p}{:}
            \PY{n}{word\PYZus{}index}\PY{o}{=}\PY{n}{word2ind}\PY{p}{[}\PY{n}{word}\PY{p}{]}
            \PY{n}{start\PYZus{}index} \PY{o}{=} \PY{n+nb}{max}\PY{p}{(}\PY{l+m+mi}{0}\PY{p}{,} \PY{n}{i} \PY{o}{\PYZhy{}} \PY{n}{window\PYZus{}size}\PY{p}{)}
            \PY{n}{end\PYZus{}index} \PY{o}{=} \PY{n+nb}{min}\PY{p}{(}\PY{n+nb}{len}\PY{p}{(}\PY{n}{doc}\PY{p}{)}\PY{p}{,} \PY{n}{i} \PY{o}{+} \PY{n}{window\PYZus{}size} \PY{o}{+} \PY{l+m+mi}{1}\PY{p}{)}
            \PY{k}{for} \PY{n}{j} \PY{o+ow}{in} \PY{n+nb}{range}\PY{p}{(}\PY{n}{start\PYZus{}index}\PY{p}{,} \PY{n}{end\PYZus{}index}\PY{p}{)}\PY{p}{:}
                \PY{k}{if} \PY{n}{j} \PY{o}{!=} \PY{n}{i}\PY{p}{:}
                    \PY{n}{co\PYZus{}word} \PY{o}{=} \PY{n}{doc}\PY{p}{[}\PY{n}{j}\PY{p}{]}
                    \PY{n}{co\PYZus{}word\PYZus{}index} \PY{o}{=} \PY{n}{word2ind}\PY{p}{[}\PY{n}{co\PYZus{}word}\PY{p}{]}
                    \PY{n}{M}\PY{p}{[}\PY{n}{word\PYZus{}index}\PY{p}{]}\PY{p}{[}\PY{n}{co\PYZus{}word\PYZus{}index}\PY{p}{]} \PY{o}{+}\PY{o}{=} \PY{l+m+mi}{1}

    \PY{c+c1}{\PYZsh{} \PYZhy{}\PYZhy{}\PYZhy{}\PYZhy{}\PYZhy{}\PYZhy{}\PYZhy{}\PYZhy{}\PYZhy{}\PYZhy{}\PYZhy{}\PYZhy{}\PYZhy{}\PYZhy{}\PYZhy{}\PYZhy{}\PYZhy{}\PYZhy{}}

    \PY{k}{return} \PY{n}{M}\PY{p}{,} \PY{n}{word2ind}
\end{Verbatim}
\end{tcolorbox}

    \begin{tcolorbox}[breakable, size=fbox, boxrule=1pt, pad at break*=1mm,colback=cellbackground, colframe=cellborder]
\prompt{In}{incolor}{ }{\boxspacing}
\begin{Verbatim}[commandchars=\\\{\}]
\PY{c+c1}{\PYZsh{} \PYZhy{}\PYZhy{}\PYZhy{}\PYZhy{}\PYZhy{}\PYZhy{}\PYZhy{}\PYZhy{}\PYZhy{}\PYZhy{}\PYZhy{}\PYZhy{}\PYZhy{}\PYZhy{}\PYZhy{}\PYZhy{}\PYZhy{}\PYZhy{}\PYZhy{}\PYZhy{}\PYZhy{}}
\PY{c+c1}{\PYZsh{} Run this sanity check}
\PY{c+c1}{\PYZsh{} Note that this is not an exhaustive check for correctness.}
\PY{c+c1}{\PYZsh{} \PYZhy{}\PYZhy{}\PYZhy{}\PYZhy{}\PYZhy{}\PYZhy{}\PYZhy{}\PYZhy{}\PYZhy{}\PYZhy{}\PYZhy{}\PYZhy{}\PYZhy{}\PYZhy{}\PYZhy{}\PYZhy{}\PYZhy{}\PYZhy{}\PYZhy{}\PYZhy{}\PYZhy{}}

\PY{c+c1}{\PYZsh{} Define toy corpus and get student\PYZsq{}s co\PYZhy{}occurrence matrix}
\PY{n}{test\PYZus{}corpus} \PY{o}{=} \PY{p}{[}\PY{l+s+s2}{\PYZdq{}}\PY{l+s+si}{\PYZob{}\PYZcb{}}\PY{l+s+s2}{ All that glitters isn}\PY{l+s+s2}{\PYZsq{}}\PY{l+s+s2}{t gold }\PY{l+s+si}{\PYZob{}\PYZcb{}}\PY{l+s+s2}{\PYZdq{}}\PY{o}{.}\PY{n}{format}\PY{p}{(}\PY{n}{START\PYZus{}TOKEN}\PY{p}{,} \PY{n}{END\PYZus{}TOKEN}\PY{p}{)}\PY{o}{.}\PY{n}{split}\PY{p}{(}\PY{l+s+s2}{\PYZdq{}}\PY{l+s+s2}{ }\PY{l+s+s2}{\PYZdq{}}\PY{p}{)}\PY{p}{,} \PY{l+s+s2}{\PYZdq{}}\PY{l+s+si}{\PYZob{}\PYZcb{}}\PY{l+s+s2}{ All}\PY{l+s+s2}{\PYZsq{}}\PY{l+s+s2}{s well that ends well }\PY{l+s+si}{\PYZob{}\PYZcb{}}\PY{l+s+s2}{\PYZdq{}}\PY{o}{.}\PY{n}{format}\PY{p}{(}\PY{n}{START\PYZus{}TOKEN}\PY{p}{,} \PY{n}{END\PYZus{}TOKEN}\PY{p}{)}\PY{o}{.}\PY{n}{split}\PY{p}{(}\PY{l+s+s2}{\PYZdq{}}\PY{l+s+s2}{ }\PY{l+s+s2}{\PYZdq{}}\PY{p}{)}\PY{p}{]}
\PY{n}{M\PYZus{}test}\PY{p}{,} \PY{n}{word2ind\PYZus{}test} \PY{o}{=} \PY{n}{compute\PYZus{}co\PYZus{}occurrence\PYZus{}matrix}\PY{p}{(}\PY{n}{test\PYZus{}corpus}\PY{p}{,} \PY{n}{window\PYZus{}size}\PY{o}{=}\PY{l+m+mi}{1}\PY{p}{)}

\PY{c+c1}{\PYZsh{} Correct M and word2ind}
\PY{n}{M\PYZus{}test\PYZus{}ans} \PY{o}{=} \PY{n}{np}\PY{o}{.}\PY{n}{array}\PY{p}{(}
    \PY{p}{[}\PY{p}{[}\PY{l+m+mf}{0.}\PY{p}{,} \PY{l+m+mf}{0.}\PY{p}{,} \PY{l+m+mf}{0.}\PY{p}{,} \PY{l+m+mf}{0.}\PY{p}{,} \PY{l+m+mf}{0.}\PY{p}{,} \PY{l+m+mf}{0.}\PY{p}{,} \PY{l+m+mf}{1.}\PY{p}{,} \PY{l+m+mf}{0.}\PY{p}{,} \PY{l+m+mf}{0.}\PY{p}{,} \PY{l+m+mf}{1.}\PY{p}{,}\PY{p}{]}\PY{p}{,}
     \PY{p}{[}\PY{l+m+mf}{0.}\PY{p}{,} \PY{l+m+mf}{0.}\PY{p}{,} \PY{l+m+mf}{1.}\PY{p}{,} \PY{l+m+mf}{1.}\PY{p}{,} \PY{l+m+mf}{0.}\PY{p}{,} \PY{l+m+mf}{0.}\PY{p}{,} \PY{l+m+mf}{0.}\PY{p}{,} \PY{l+m+mf}{0.}\PY{p}{,} \PY{l+m+mf}{0.}\PY{p}{,} \PY{l+m+mf}{0.}\PY{p}{,}\PY{p}{]}\PY{p}{,}
     \PY{p}{[}\PY{l+m+mf}{0.}\PY{p}{,} \PY{l+m+mf}{1.}\PY{p}{,} \PY{l+m+mf}{0.}\PY{p}{,} \PY{l+m+mf}{0.}\PY{p}{,} \PY{l+m+mf}{0.}\PY{p}{,} \PY{l+m+mf}{0.}\PY{p}{,} \PY{l+m+mf}{0.}\PY{p}{,} \PY{l+m+mf}{0.}\PY{p}{,} \PY{l+m+mf}{1.}\PY{p}{,} \PY{l+m+mf}{0.}\PY{p}{,}\PY{p}{]}\PY{p}{,}
     \PY{p}{[}\PY{l+m+mf}{0.}\PY{p}{,} \PY{l+m+mf}{1.}\PY{p}{,} \PY{l+m+mf}{0.}\PY{p}{,} \PY{l+m+mf}{0.}\PY{p}{,} \PY{l+m+mf}{0.}\PY{p}{,} \PY{l+m+mf}{0.}\PY{p}{,} \PY{l+m+mf}{0.}\PY{p}{,} \PY{l+m+mf}{0.}\PY{p}{,} \PY{l+m+mf}{0.}\PY{p}{,} \PY{l+m+mf}{1.}\PY{p}{,}\PY{p}{]}\PY{p}{,}
     \PY{p}{[}\PY{l+m+mf}{0.}\PY{p}{,} \PY{l+m+mf}{0.}\PY{p}{,} \PY{l+m+mf}{0.}\PY{p}{,} \PY{l+m+mf}{0.}\PY{p}{,} \PY{l+m+mf}{0.}\PY{p}{,} \PY{l+m+mf}{0.}\PY{p}{,} \PY{l+m+mf}{0.}\PY{p}{,} \PY{l+m+mf}{0.}\PY{p}{,} \PY{l+m+mf}{1.}\PY{p}{,} \PY{l+m+mf}{1.}\PY{p}{,}\PY{p}{]}\PY{p}{,}
     \PY{p}{[}\PY{l+m+mf}{0.}\PY{p}{,} \PY{l+m+mf}{0.}\PY{p}{,} \PY{l+m+mf}{0.}\PY{p}{,} \PY{l+m+mf}{0.}\PY{p}{,} \PY{l+m+mf}{0.}\PY{p}{,} \PY{l+m+mf}{0.}\PY{p}{,} \PY{l+m+mf}{0.}\PY{p}{,} \PY{l+m+mf}{1.}\PY{p}{,} \PY{l+m+mf}{1.}\PY{p}{,} \PY{l+m+mf}{0.}\PY{p}{,}\PY{p}{]}\PY{p}{,}
     \PY{p}{[}\PY{l+m+mf}{1.}\PY{p}{,} \PY{l+m+mf}{0.}\PY{p}{,} \PY{l+m+mf}{0.}\PY{p}{,} \PY{l+m+mf}{0.}\PY{p}{,} \PY{l+m+mf}{0.}\PY{p}{,} \PY{l+m+mf}{0.}\PY{p}{,} \PY{l+m+mf}{0.}\PY{p}{,} \PY{l+m+mf}{1.}\PY{p}{,} \PY{l+m+mf}{0.}\PY{p}{,} \PY{l+m+mf}{0.}\PY{p}{,}\PY{p}{]}\PY{p}{,}
     \PY{p}{[}\PY{l+m+mf}{0.}\PY{p}{,} \PY{l+m+mf}{0.}\PY{p}{,} \PY{l+m+mf}{0.}\PY{p}{,} \PY{l+m+mf}{0.}\PY{p}{,} \PY{l+m+mf}{0.}\PY{p}{,} \PY{l+m+mf}{1.}\PY{p}{,} \PY{l+m+mf}{1.}\PY{p}{,} \PY{l+m+mf}{0.}\PY{p}{,} \PY{l+m+mf}{0.}\PY{p}{,} \PY{l+m+mf}{0.}\PY{p}{,}\PY{p}{]}\PY{p}{,}
     \PY{p}{[}\PY{l+m+mf}{0.}\PY{p}{,} \PY{l+m+mf}{0.}\PY{p}{,} \PY{l+m+mf}{1.}\PY{p}{,} \PY{l+m+mf}{0.}\PY{p}{,} \PY{l+m+mf}{1.}\PY{p}{,} \PY{l+m+mf}{1.}\PY{p}{,} \PY{l+m+mf}{0.}\PY{p}{,} \PY{l+m+mf}{0.}\PY{p}{,} \PY{l+m+mf}{0.}\PY{p}{,} \PY{l+m+mf}{1.}\PY{p}{,}\PY{p}{]}\PY{p}{,}
     \PY{p}{[}\PY{l+m+mf}{1.}\PY{p}{,} \PY{l+m+mf}{0.}\PY{p}{,} \PY{l+m+mf}{0.}\PY{p}{,} \PY{l+m+mf}{1.}\PY{p}{,} \PY{l+m+mf}{1.}\PY{p}{,} \PY{l+m+mf}{0.}\PY{p}{,} \PY{l+m+mf}{0.}\PY{p}{,} \PY{l+m+mf}{0.}\PY{p}{,} \PY{l+m+mf}{1.}\PY{p}{,} \PY{l+m+mf}{0.}\PY{p}{,}\PY{p}{]}\PY{p}{]}
\PY{p}{)}
\PY{n}{ans\PYZus{}test\PYZus{}corpus\PYZus{}words} \PY{o}{=} \PY{n+nb}{sorted}\PY{p}{(}\PY{p}{[}\PY{n}{START\PYZus{}TOKEN}\PY{p}{,} \PY{l+s+s2}{\PYZdq{}}\PY{l+s+s2}{All}\PY{l+s+s2}{\PYZdq{}}\PY{p}{,} \PY{l+s+s2}{\PYZdq{}}\PY{l+s+s2}{ends}\PY{l+s+s2}{\PYZdq{}}\PY{p}{,} \PY{l+s+s2}{\PYZdq{}}\PY{l+s+s2}{that}\PY{l+s+s2}{\PYZdq{}}\PY{p}{,} \PY{l+s+s2}{\PYZdq{}}\PY{l+s+s2}{gold}\PY{l+s+s2}{\PYZdq{}}\PY{p}{,} \PY{l+s+s2}{\PYZdq{}}\PY{l+s+s2}{All}\PY{l+s+s2}{\PYZsq{}}\PY{l+s+s2}{s}\PY{l+s+s2}{\PYZdq{}}\PY{p}{,} \PY{l+s+s2}{\PYZdq{}}\PY{l+s+s2}{glitters}\PY{l+s+s2}{\PYZdq{}}\PY{p}{,} \PY{l+s+s2}{\PYZdq{}}\PY{l+s+s2}{isn}\PY{l+s+s2}{\PYZsq{}}\PY{l+s+s2}{t}\PY{l+s+s2}{\PYZdq{}}\PY{p}{,} \PY{l+s+s2}{\PYZdq{}}\PY{l+s+s2}{well}\PY{l+s+s2}{\PYZdq{}}\PY{p}{,} \PY{n}{END\PYZus{}TOKEN}\PY{p}{]}\PY{p}{)}
\PY{n}{word2ind\PYZus{}ans} \PY{o}{=} \PY{n+nb}{dict}\PY{p}{(}\PY{n+nb}{zip}\PY{p}{(}\PY{n}{ans\PYZus{}test\PYZus{}corpus\PYZus{}words}\PY{p}{,} \PY{n+nb}{range}\PY{p}{(}\PY{n+nb}{len}\PY{p}{(}\PY{n}{ans\PYZus{}test\PYZus{}corpus\PYZus{}words}\PY{p}{)}\PY{p}{)}\PY{p}{)}\PY{p}{)}

\PY{c+c1}{\PYZsh{} Test correct word2ind}
\PY{k}{assert} \PY{p}{(}\PY{n}{word2ind\PYZus{}ans} \PY{o}{==} \PY{n}{word2ind\PYZus{}test}\PY{p}{)}\PY{p}{,} \PY{l+s+s2}{\PYZdq{}}\PY{l+s+s2}{Your word2ind is incorrect:}\PY{l+s+se}{\PYZbs{}n}\PY{l+s+s2}{Correct: }\PY{l+s+si}{\PYZob{}\PYZcb{}}\PY{l+s+se}{\PYZbs{}n}\PY{l+s+s2}{Yours: }\PY{l+s+si}{\PYZob{}\PYZcb{}}\PY{l+s+s2}{\PYZdq{}}\PY{o}{.}\PY{n}{format}\PY{p}{(}\PY{n}{word2ind\PYZus{}ans}\PY{p}{,} \PY{n}{word2ind\PYZus{}test}\PY{p}{)}

\PY{c+c1}{\PYZsh{} Test correct M shape}
\PY{k}{assert} \PY{p}{(}\PY{n}{M\PYZus{}test}\PY{o}{.}\PY{n}{shape} \PY{o}{==} \PY{n}{M\PYZus{}test\PYZus{}ans}\PY{o}{.}\PY{n}{shape}\PY{p}{)}\PY{p}{,} \PY{l+s+s2}{\PYZdq{}}\PY{l+s+s2}{M matrix has incorrect shape.}\PY{l+s+se}{\PYZbs{}n}\PY{l+s+s2}{Correct: }\PY{l+s+si}{\PYZob{}\PYZcb{}}\PY{l+s+se}{\PYZbs{}n}\PY{l+s+s2}{Yours: }\PY{l+s+si}{\PYZob{}\PYZcb{}}\PY{l+s+s2}{\PYZdq{}}\PY{o}{.}\PY{n}{format}\PY{p}{(}\PY{n}{M\PYZus{}test}\PY{o}{.}\PY{n}{shape}\PY{p}{,} \PY{n}{M\PYZus{}test\PYZus{}ans}\PY{o}{.}\PY{n}{shape}\PY{p}{)}

\PY{c+c1}{\PYZsh{} Test correct M values}
\PY{k}{for} \PY{n}{w1} \PY{o+ow}{in} \PY{n}{word2ind\PYZus{}ans}\PY{o}{.}\PY{n}{keys}\PY{p}{(}\PY{p}{)}\PY{p}{:}
    \PY{n}{idx1} \PY{o}{=} \PY{n}{word2ind\PYZus{}ans}\PY{p}{[}\PY{n}{w1}\PY{p}{]}
    \PY{k}{for} \PY{n}{w2} \PY{o+ow}{in} \PY{n}{word2ind\PYZus{}ans}\PY{o}{.}\PY{n}{keys}\PY{p}{(}\PY{p}{)}\PY{p}{:}
        \PY{n}{idx2} \PY{o}{=} \PY{n}{word2ind\PYZus{}ans}\PY{p}{[}\PY{n}{w2}\PY{p}{]}
        \PY{n}{student} \PY{o}{=} \PY{n}{M\PYZus{}test}\PY{p}{[}\PY{n}{idx1}\PY{p}{,} \PY{n}{idx2}\PY{p}{]}
        \PY{n}{correct} \PY{o}{=} \PY{n}{M\PYZus{}test\PYZus{}ans}\PY{p}{[}\PY{n}{idx1}\PY{p}{,} \PY{n}{idx2}\PY{p}{]}
        \PY{k}{if} \PY{n}{student} \PY{o}{!=} \PY{n}{correct}\PY{p}{:}
            \PY{n+nb}{print}\PY{p}{(}\PY{l+s+s2}{\PYZdq{}}\PY{l+s+s2}{Correct M:}\PY{l+s+s2}{\PYZdq{}}\PY{p}{)}
            \PY{n+nb}{print}\PY{p}{(}\PY{n}{M\PYZus{}test\PYZus{}ans}\PY{p}{)}
            \PY{n+nb}{print}\PY{p}{(}\PY{l+s+s2}{\PYZdq{}}\PY{l+s+s2}{Your M: }\PY{l+s+s2}{\PYZdq{}}\PY{p}{)}
            \PY{n+nb}{print}\PY{p}{(}\PY{n}{M\PYZus{}test}\PY{p}{)}
            \PY{k}{raise} \PY{n+ne}{AssertionError}\PY{p}{(}\PY{l+s+s2}{\PYZdq{}}\PY{l+s+s2}{Incorrect count at index (}\PY{l+s+si}{\PYZob{}\PYZcb{}}\PY{l+s+s2}{, }\PY{l+s+si}{\PYZob{}\PYZcb{}}\PY{l+s+s2}{)=(}\PY{l+s+si}{\PYZob{}\PYZcb{}}\PY{l+s+s2}{, }\PY{l+s+si}{\PYZob{}\PYZcb{}}\PY{l+s+s2}{) in matrix M. Yours has }\PY{l+s+si}{\PYZob{}\PYZcb{}}\PY{l+s+s2}{ but should have }\PY{l+s+si}{\PYZob{}\PYZcb{}}\PY{l+s+s2}{.}\PY{l+s+s2}{\PYZdq{}}\PY{o}{.}\PY{n}{format}\PY{p}{(}\PY{n}{idx1}\PY{p}{,} \PY{n}{idx2}\PY{p}{,} \PY{n}{w1}\PY{p}{,} \PY{n}{w2}\PY{p}{,} \PY{n}{student}\PY{p}{,} \PY{n}{correct}\PY{p}{)}\PY{p}{)}

\PY{c+c1}{\PYZsh{} Print Success}
\PY{n+nb}{print} \PY{p}{(}\PY{l+s+s2}{\PYZdq{}}\PY{l+s+s2}{\PYZhy{}}\PY{l+s+s2}{\PYZdq{}} \PY{o}{*} \PY{l+m+mi}{80}\PY{p}{)}
\PY{n+nb}{print}\PY{p}{(}\PY{l+s+s2}{\PYZdq{}}\PY{l+s+s2}{Passed All Tests!}\PY{l+s+s2}{\PYZdq{}}\PY{p}{)}
\PY{n+nb}{print} \PY{p}{(}\PY{l+s+s2}{\PYZdq{}}\PY{l+s+s2}{\PYZhy{}}\PY{l+s+s2}{\PYZdq{}} \PY{o}{*} \PY{l+m+mi}{80}\PY{p}{)}
\end{Verbatim}
\end{tcolorbox}

    \begin{Verbatim}[commandchars=\\\{\}]
--------------------------------------------------------------------------------
Passed All Tests!
--------------------------------------------------------------------------------
    \end{Verbatim}

    \subsubsection{\texorpdfstring{Question 1.3: Implement
\texttt{reduce\_to\_k\_dim} {[}code{]} (1
point)}{Question 1.3: Implement reduce\_to\_k\_dim {[}code{]} (1 point)}}\label{question-1.3-implement-reduce_to_k_dim-code-1-point}

Construct a method that performs dimensionality reduction on the matrix
to produce k-dimensional embeddings. Use SVD to take the top k
components and produce a new matrix of k-dimensional embeddings.

\textbf{Note:} All of numpy, scipy, and scikit-learn (\texttt{sklearn})
provide \emph{some} implementation of SVD, but only scipy and sklearn
provide an implementation of Truncated SVD, and only sklearn provides an
efficient randomized algorithm for calculating large-scale Truncated
SVD. So please use
\href{https://scikit-learn.org/stable/modules/generated/sklearn.decomposition.TruncatedSVD.html}{sklearn.decomposition.TruncatedSVD}.

    \begin{tcolorbox}[breakable, size=fbox, boxrule=1pt, pad at break*=1mm,colback=cellbackground, colframe=cellborder]
\prompt{In}{incolor}{ }{\boxspacing}
\begin{Verbatim}[commandchars=\\\{\}]
\PY{k}{def}\PY{+w}{ }\PY{n+nf}{reduce\PYZus{}to\PYZus{}k\PYZus{}dim}\PY{p}{(}\PY{n}{M}\PY{p}{,} \PY{n}{k}\PY{o}{=}\PY{l+m+mi}{2}\PY{p}{)}\PY{p}{:}
\PY{+w}{    }\PY{l+s+sd}{\PYZdq{}\PYZdq{}\PYZdq{} Reduce a co\PYZhy{}occurence count matrix of dimensionality (num\PYZus{}corpus\PYZus{}words, num\PYZus{}corpus\PYZus{}words)}
\PY{l+s+sd}{        to a matrix of dimensionality (num\PYZus{}corpus\PYZus{}words, k) using the following SVD function from Scikit\PYZhy{}Learn:}
\PY{l+s+sd}{            \PYZhy{} http://scikit\PYZhy{}learn.org/stable/modules/generated/sklearn.decomposition.TruncatedSVD.html}

\PY{l+s+sd}{        Params:}
\PY{l+s+sd}{            M (numpy matrix of shape (number of unique words in the corpus , number of unique words in the corpus)): co\PYZhy{}occurence matrix of word counts}
\PY{l+s+sd}{            k (int): embedding size of each word after dimension reduction}
\PY{l+s+sd}{        Return:}
\PY{l+s+sd}{            M\PYZus{}reduced (numpy matrix of shape (number of corpus words, k)): matrix of k\PYZhy{}dimensioal word embeddings.}
\PY{l+s+sd}{                    In terms of the SVD from math class, this actually returns U * S}
\PY{l+s+sd}{    \PYZdq{}\PYZdq{}\PYZdq{}}
    \PY{n}{n\PYZus{}iters} \PY{o}{=} \PY{l+m+mi}{10}    \PY{c+c1}{\PYZsh{} Use this parameter in your call to `TruncatedSVD`}
    \PY{n}{M\PYZus{}reduced} \PY{o}{=} \PY{k+kc}{None}
    \PY{n+nb}{print}\PY{p}{(}\PY{l+s+s2}{\PYZdq{}}\PY{l+s+s2}{Running Truncated SVD over }\PY{l+s+si}{\PYZpc{}i}\PY{l+s+s2}{ words...}\PY{l+s+s2}{\PYZdq{}} \PY{o}{\PYZpc{}} \PY{p}{(}\PY{n}{M}\PY{o}{.}\PY{n}{shape}\PY{p}{[}\PY{l+m+mi}{0}\PY{p}{]}\PY{p}{)}\PY{p}{)}

    \PY{c+c1}{\PYZsh{} \PYZhy{}\PYZhy{}\PYZhy{}\PYZhy{}\PYZhy{}\PYZhy{}\PYZhy{}\PYZhy{}\PYZhy{}\PYZhy{}\PYZhy{}\PYZhy{}\PYZhy{}\PYZhy{}\PYZhy{}\PYZhy{}\PYZhy{}\PYZhy{}}
    \PY{c+c1}{\PYZsh{} Write your implementation here.}
    \PY{n}{TruncatedSVD\PYZus{}model} \PY{o}{=} \PY{n}{TruncatedSVD}\PY{p}{(}\PY{n}{n\PYZus{}components}\PY{o}{=}\PY{n}{k}\PY{p}{,} \PY{n}{n\PYZus{}iter}\PY{o}{=}\PY{n}{n\PYZus{}iters}\PY{p}{)}
    \PY{n}{M\PYZus{}reduced} \PY{o}{=} \PY{n}{TruncatedSVD\PYZus{}model}\PY{o}{.}\PY{n}{fit\PYZus{}transform}\PY{p}{(}\PY{n}{M}\PY{p}{)}


    \PY{c+c1}{\PYZsh{} \PYZhy{}\PYZhy{}\PYZhy{}\PYZhy{}\PYZhy{}\PYZhy{}\PYZhy{}\PYZhy{}\PYZhy{}\PYZhy{}\PYZhy{}\PYZhy{}\PYZhy{}\PYZhy{}\PYZhy{}\PYZhy{}\PYZhy{}\PYZhy{}}

    \PY{n+nb}{print}\PY{p}{(}\PY{l+s+s2}{\PYZdq{}}\PY{l+s+s2}{Done.}\PY{l+s+s2}{\PYZdq{}}\PY{p}{)}
    \PY{k}{return} \PY{n}{M\PYZus{}reduced}
\end{Verbatim}
\end{tcolorbox}

    \begin{tcolorbox}[breakable, size=fbox, boxrule=1pt, pad at break*=1mm,colback=cellbackground, colframe=cellborder]
\prompt{In}{incolor}{ }{\boxspacing}
\begin{Verbatim}[commandchars=\\\{\}]
\PY{c+c1}{\PYZsh{} \PYZhy{}\PYZhy{}\PYZhy{}\PYZhy{}\PYZhy{}\PYZhy{}\PYZhy{}\PYZhy{}\PYZhy{}\PYZhy{}\PYZhy{}\PYZhy{}\PYZhy{}\PYZhy{}\PYZhy{}\PYZhy{}\PYZhy{}\PYZhy{}\PYZhy{}\PYZhy{}\PYZhy{}}
\PY{c+c1}{\PYZsh{} Run this sanity check}
\PY{c+c1}{\PYZsh{} Note that this is not an exhaustive check for correctness}
\PY{c+c1}{\PYZsh{} In fact we only check that your M\PYZus{}reduced has the right dimensions.}
\PY{c+c1}{\PYZsh{} \PYZhy{}\PYZhy{}\PYZhy{}\PYZhy{}\PYZhy{}\PYZhy{}\PYZhy{}\PYZhy{}\PYZhy{}\PYZhy{}\PYZhy{}\PYZhy{}\PYZhy{}\PYZhy{}\PYZhy{}\PYZhy{}\PYZhy{}\PYZhy{}\PYZhy{}\PYZhy{}\PYZhy{}}

\PY{c+c1}{\PYZsh{} Define toy corpus and run student code}
\PY{n}{test\PYZus{}corpus} \PY{o}{=} \PY{p}{[}\PY{l+s+s2}{\PYZdq{}}\PY{l+s+si}{\PYZob{}\PYZcb{}}\PY{l+s+s2}{ All that glitters isn}\PY{l+s+s2}{\PYZsq{}}\PY{l+s+s2}{t gold }\PY{l+s+si}{\PYZob{}\PYZcb{}}\PY{l+s+s2}{\PYZdq{}}\PY{o}{.}\PY{n}{format}\PY{p}{(}\PY{n}{START\PYZus{}TOKEN}\PY{p}{,} \PY{n}{END\PYZus{}TOKEN}\PY{p}{)}\PY{o}{.}\PY{n}{split}\PY{p}{(}\PY{l+s+s2}{\PYZdq{}}\PY{l+s+s2}{ }\PY{l+s+s2}{\PYZdq{}}\PY{p}{)}\PY{p}{,} \PY{l+s+s2}{\PYZdq{}}\PY{l+s+si}{\PYZob{}\PYZcb{}}\PY{l+s+s2}{ All}\PY{l+s+s2}{\PYZsq{}}\PY{l+s+s2}{s well that ends well }\PY{l+s+si}{\PYZob{}\PYZcb{}}\PY{l+s+s2}{\PYZdq{}}\PY{o}{.}\PY{n}{format}\PY{p}{(}\PY{n}{START\PYZus{}TOKEN}\PY{p}{,} \PY{n}{END\PYZus{}TOKEN}\PY{p}{)}\PY{o}{.}\PY{n}{split}\PY{p}{(}\PY{l+s+s2}{\PYZdq{}}\PY{l+s+s2}{ }\PY{l+s+s2}{\PYZdq{}}\PY{p}{)}\PY{p}{]}
\PY{n}{M\PYZus{}test}\PY{p}{,} \PY{n}{word2ind\PYZus{}test} \PY{o}{=} \PY{n}{compute\PYZus{}co\PYZus{}occurrence\PYZus{}matrix}\PY{p}{(}\PY{n}{test\PYZus{}corpus}\PY{p}{,} \PY{n}{window\PYZus{}size}\PY{o}{=}\PY{l+m+mi}{1}\PY{p}{)}
\PY{n}{M\PYZus{}test\PYZus{}reduced} \PY{o}{=} \PY{n}{reduce\PYZus{}to\PYZus{}k\PYZus{}dim}\PY{p}{(}\PY{n}{M\PYZus{}test}\PY{p}{,} \PY{n}{k}\PY{o}{=}\PY{l+m+mi}{2}\PY{p}{)}

\PY{c+c1}{\PYZsh{} Test proper dimensions}
\PY{k}{assert} \PY{p}{(}\PY{n}{M\PYZus{}test\PYZus{}reduced}\PY{o}{.}\PY{n}{shape}\PY{p}{[}\PY{l+m+mi}{0}\PY{p}{]} \PY{o}{==} \PY{l+m+mi}{10}\PY{p}{)}\PY{p}{,} \PY{l+s+s2}{\PYZdq{}}\PY{l+s+s2}{M\PYZus{}reduced has }\PY{l+s+si}{\PYZob{}\PYZcb{}}\PY{l+s+s2}{ rows; should have }\PY{l+s+si}{\PYZob{}\PYZcb{}}\PY{l+s+s2}{\PYZdq{}}\PY{o}{.}\PY{n}{format}\PY{p}{(}\PY{n}{M\PYZus{}test\PYZus{}reduced}\PY{o}{.}\PY{n}{shape}\PY{p}{[}\PY{l+m+mi}{0}\PY{p}{]}\PY{p}{,} \PY{l+m+mi}{10}\PY{p}{)}
\PY{k}{assert} \PY{p}{(}\PY{n}{M\PYZus{}test\PYZus{}reduced}\PY{o}{.}\PY{n}{shape}\PY{p}{[}\PY{l+m+mi}{1}\PY{p}{]} \PY{o}{==} \PY{l+m+mi}{2}\PY{p}{)}\PY{p}{,} \PY{l+s+s2}{\PYZdq{}}\PY{l+s+s2}{M\PYZus{}reduced has }\PY{l+s+si}{\PYZob{}\PYZcb{}}\PY{l+s+s2}{ columns; should have }\PY{l+s+si}{\PYZob{}\PYZcb{}}\PY{l+s+s2}{\PYZdq{}}\PY{o}{.}\PY{n}{format}\PY{p}{(}\PY{n}{M\PYZus{}test\PYZus{}reduced}\PY{o}{.}\PY{n}{shape}\PY{p}{[}\PY{l+m+mi}{1}\PY{p}{]}\PY{p}{,} \PY{l+m+mi}{2}\PY{p}{)}

\PY{c+c1}{\PYZsh{} Print Success}
\PY{n+nb}{print} \PY{p}{(}\PY{l+s+s2}{\PYZdq{}}\PY{l+s+s2}{\PYZhy{}}\PY{l+s+s2}{\PYZdq{}} \PY{o}{*} \PY{l+m+mi}{80}\PY{p}{)}
\PY{n+nb}{print}\PY{p}{(}\PY{l+s+s2}{\PYZdq{}}\PY{l+s+s2}{Passed All Tests!}\PY{l+s+s2}{\PYZdq{}}\PY{p}{)}
\PY{n+nb}{print} \PY{p}{(}\PY{l+s+s2}{\PYZdq{}}\PY{l+s+s2}{\PYZhy{}}\PY{l+s+s2}{\PYZdq{}} \PY{o}{*} \PY{l+m+mi}{80}\PY{p}{)}
\end{Verbatim}
\end{tcolorbox}

    \begin{Verbatim}[commandchars=\\\{\}]
Running Truncated SVD over 10 words{\ldots}
Done.
--------------------------------------------------------------------------------
Passed All Tests!
--------------------------------------------------------------------------------
    \end{Verbatim}

    \subsubsection{\texorpdfstring{Question 1.4: Implement
\texttt{plot\_embeddings} {[}code{]} (1
point)}{Question 1.4: Implement plot\_embeddings {[}code{]} (1 point)}}\label{question-1.4-implement-plot_embeddings-code-1-point}

Here you will write a function to plot a set of 2D vectors in 2D space.
For graphs, we will use Matplotlib (\texttt{plt}).

For this example, you may find it useful to adapt
\href{http://web.archive.org/web/20190924160434/https://www.pythonmembers.club/2018/05/08/matplotlib-scatter-plot-annotate-set-text-at-label-each-point/}{this
code}. In the future, a good way to make a plot is to look at
\href{https://matplotlib.org/gallery/index.html}{the Matplotlib
gallery}, find a plot that looks somewhat like what you want, and adapt
the code they give.

    \begin{tcolorbox}[breakable, size=fbox, boxrule=1pt, pad at break*=1mm,colback=cellbackground, colframe=cellborder]
\prompt{In}{incolor}{ }{\boxspacing}
\begin{Verbatim}[commandchars=\\\{\}]
\PY{k}{def}\PY{+w}{ }\PY{n+nf}{plot\PYZus{}embeddings}\PY{p}{(}\PY{n}{M\PYZus{}reduced}\PY{p}{,} \PY{n}{word2ind}\PY{p}{,} \PY{n}{words}\PY{p}{)}\PY{p}{:}
\PY{+w}{    }\PY{l+s+sd}{\PYZdq{}\PYZdq{}\PYZdq{} Plot in a scatterplot the embeddings of the words specified in the list \PYZdq{}words\PYZdq{}.}
\PY{l+s+sd}{        NOTE: do not plot all the words listed in M\PYZus{}reduced / word2ind.}
\PY{l+s+sd}{        Include a label next to each point.}

\PY{l+s+sd}{        Params:}
\PY{l+s+sd}{            M\PYZus{}reduced (numpy matrix of shape (number of unique words in the corpus , 2)): matrix of 2\PYZhy{}dimensioal word embeddings}
\PY{l+s+sd}{            word2ind (dict): dictionary that maps word to indices for matrix M}
\PY{l+s+sd}{            words (list of strings): words whose embeddings we want to visualize}
\PY{l+s+sd}{    \PYZdq{}\PYZdq{}\PYZdq{}}

    \PY{c+c1}{\PYZsh{} \PYZhy{}\PYZhy{}\PYZhy{}\PYZhy{}\PYZhy{}\PYZhy{}\PYZhy{}\PYZhy{}\PYZhy{}\PYZhy{}\PYZhy{}\PYZhy{}\PYZhy{}\PYZhy{}\PYZhy{}\PYZhy{}\PYZhy{}\PYZhy{}}
    \PY{c+c1}{\PYZsh{} Write your implementation here.}
    \PY{n}{x\PYZus{}coords} \PY{o}{=} \PY{p}{[}\PY{p}{]}
    \PY{n}{y\PYZus{}coords} \PY{o}{=} \PY{p}{[}\PY{p}{]}
    \PY{n}{labels} \PY{o}{=} \PY{p}{[}\PY{p}{]}

    \PY{k}{for} \PY{n}{word} \PY{o+ow}{in} \PY{n}{words}\PY{p}{:}
        \PY{k}{if} \PY{n}{word} \PY{o+ow}{in} \PY{n}{word2ind}\PY{p}{:}
            \PY{n}{idx} \PY{o}{=} \PY{n}{word2ind}\PY{p}{[}\PY{n}{word}\PY{p}{]}
            \PY{n}{x\PYZus{}coords}\PY{o}{.}\PY{n}{append}\PY{p}{(}\PY{n}{M\PYZus{}reduced}\PY{p}{[}\PY{n}{idx}\PY{p}{,} \PY{l+m+mi}{0}\PY{p}{]}\PY{p}{)}
            \PY{n}{y\PYZus{}coords}\PY{o}{.}\PY{n}{append}\PY{p}{(}\PY{n}{M\PYZus{}reduced}\PY{p}{[}\PY{n}{idx}\PY{p}{,} \PY{l+m+mi}{1}\PY{p}{]}\PY{p}{)}
            \PY{n}{labels}\PY{o}{.}\PY{n}{append}\PY{p}{(}\PY{n}{word}\PY{p}{)}

    \PY{c+c1}{\PYZsh{} 绘制散点图}
    \PY{n}{plt}\PY{o}{.}\PY{n}{figure}\PY{p}{(}\PY{n}{figsize}\PY{o}{=}\PY{p}{(}\PY{l+m+mi}{10}\PY{p}{,} \PY{l+m+mi}{8}\PY{p}{)}\PY{p}{)}
    \PY{n}{plt}\PY{o}{.}\PY{n}{scatter}\PY{p}{(}\PY{n}{x\PYZus{}coords}\PY{p}{,} \PY{n}{y\PYZus{}coords}\PY{p}{,} \PY{n}{marker}\PY{o}{=}\PY{l+s+s1}{\PYZsq{}}\PY{l+s+s1}{o}\PY{l+s+s1}{\PYZsq{}}\PY{p}{,} \PY{n}{color}\PY{o}{=}\PY{l+s+s1}{\PYZsq{}}\PY{l+s+s1}{blue}\PY{l+s+s1}{\PYZsq{}}\PY{p}{,} \PY{n}{s}\PY{o}{=}\PY{l+m+mi}{50}\PY{p}{)}

    \PY{c+c1}{\PYZsh{} 为每个点添加文字标签}
    \PY{k}{for} \PY{n}{x}\PY{p}{,} \PY{n}{y}\PY{p}{,} \PY{n}{label} \PY{o+ow}{in} \PY{n+nb}{zip}\PY{p}{(}\PY{n}{x\PYZus{}coords}\PY{p}{,} \PY{n}{y\PYZus{}coords}\PY{p}{,} \PY{n}{labels}\PY{p}{)}\PY{p}{:}
        \PY{n}{plt}\PY{o}{.}\PY{n}{annotate}\PY{p}{(}\PY{n}{label}\PY{p}{,} \PY{n}{xy}\PY{o}{=}\PY{p}{(}\PY{n}{x}\PY{p}{,} \PY{n}{y}\PY{p}{)}\PY{p}{,} \PY{n}{xytext}\PY{o}{=}\PY{p}{(}\PY{l+m+mi}{5}\PY{p}{,} \PY{l+m+mi}{2}\PY{p}{)}\PY{p}{,}
                     \PY{n}{textcoords}\PY{o}{=}\PY{l+s+s1}{\PYZsq{}}\PY{l+s+s1}{offset points}\PY{l+s+s1}{\PYZsq{}}\PY{p}{,} \PY{n}{fontsize}\PY{o}{=}\PY{l+m+mi}{12}\PY{p}{,} \PY{n}{color}\PY{o}{=}\PY{l+s+s1}{\PYZsq{}}\PY{l+s+s1}{darkred}\PY{l+s+s1}{\PYZsq{}}\PY{p}{)}

    \PY{n}{plt}\PY{o}{.}\PY{n}{xlabel}\PY{p}{(}\PY{l+s+s1}{\PYZsq{}}\PY{l+s+s1}{Dimension 1}\PY{l+s+s1}{\PYZsq{}}\PY{p}{)}
    \PY{n}{plt}\PY{o}{.}\PY{n}{ylabel}\PY{p}{(}\PY{l+s+s1}{\PYZsq{}}\PY{l+s+s1}{Dimension 2}\PY{l+s+s1}{\PYZsq{}}\PY{p}{)}
    \PY{n}{plt}\PY{o}{.}\PY{n}{title}\PY{p}{(}\PY{l+s+s1}{\PYZsq{}}\PY{l+s+s1}{Word Embeddings}\PY{l+s+s1}{\PYZsq{}}\PY{p}{)}
    \PY{n}{plt}\PY{o}{.}\PY{n}{grid}\PY{p}{(}\PY{k+kc}{True}\PY{p}{,} \PY{n}{linestyle}\PY{o}{=}\PY{l+s+s1}{\PYZsq{}}\PY{l+s+s1}{\PYZhy{}\PYZhy{}}\PY{l+s+s1}{\PYZsq{}}\PY{p}{,} \PY{n}{alpha}\PY{o}{=}\PY{l+m+mf}{0.5}\PY{p}{)}
    \PY{n}{plt}\PY{o}{.}\PY{n}{show}\PY{p}{(}\PY{p}{)}

    \PY{c+c1}{\PYZsh{} \PYZhy{}\PYZhy{}\PYZhy{}\PYZhy{}\PYZhy{}\PYZhy{}\PYZhy{}\PYZhy{}\PYZhy{}\PYZhy{}\PYZhy{}\PYZhy{}\PYZhy{}\PYZhy{}\PYZhy{}\PYZhy{}\PYZhy{}\PYZhy{}}
\end{Verbatim}
\end{tcolorbox}

    \begin{tcolorbox}[breakable, size=fbox, boxrule=1pt, pad at break*=1mm,colback=cellbackground, colframe=cellborder]
\prompt{In}{incolor}{ }{\boxspacing}
\begin{Verbatim}[commandchars=\\\{\}]
\PY{c+c1}{\PYZsh{} \PYZhy{}\PYZhy{}\PYZhy{}\PYZhy{}\PYZhy{}\PYZhy{}\PYZhy{}\PYZhy{}\PYZhy{}\PYZhy{}\PYZhy{}\PYZhy{}\PYZhy{}\PYZhy{}\PYZhy{}\PYZhy{}\PYZhy{}\PYZhy{}\PYZhy{}\PYZhy{}\PYZhy{}}
\PY{c+c1}{\PYZsh{} Run this sanity check}
\PY{c+c1}{\PYZsh{} Note that this is not an exhaustive check for correctness.}
\PY{c+c1}{\PYZsh{} The plot produced should look like the included file question\PYZus{}1.4\PYZus{}test.png}
\PY{c+c1}{\PYZsh{} \PYZhy{}\PYZhy{}\PYZhy{}\PYZhy{}\PYZhy{}\PYZhy{}\PYZhy{}\PYZhy{}\PYZhy{}\PYZhy{}\PYZhy{}\PYZhy{}\PYZhy{}\PYZhy{}\PYZhy{}\PYZhy{}\PYZhy{}\PYZhy{}\PYZhy{}\PYZhy{}\PYZhy{}}

\PY{n+nb}{print} \PY{p}{(}\PY{l+s+s2}{\PYZdq{}}\PY{l+s+s2}{\PYZhy{}}\PY{l+s+s2}{\PYZdq{}} \PY{o}{*} \PY{l+m+mi}{80}\PY{p}{)}
\PY{n+nb}{print} \PY{p}{(}\PY{l+s+s2}{\PYZdq{}}\PY{l+s+s2}{Outputted Plot:}\PY{l+s+s2}{\PYZdq{}}\PY{p}{)}

\PY{n}{M\PYZus{}reduced\PYZus{}plot\PYZus{}test} \PY{o}{=} \PY{n}{np}\PY{o}{.}\PY{n}{array}\PY{p}{(}\PY{p}{[}\PY{p}{[}\PY{l+m+mi}{1}\PY{p}{,} \PY{l+m+mi}{1}\PY{p}{]}\PY{p}{,} \PY{p}{[}\PY{o}{\PYZhy{}}\PY{l+m+mi}{1}\PY{p}{,} \PY{o}{\PYZhy{}}\PY{l+m+mi}{1}\PY{p}{]}\PY{p}{,} \PY{p}{[}\PY{l+m+mi}{1}\PY{p}{,} \PY{o}{\PYZhy{}}\PY{l+m+mi}{1}\PY{p}{]}\PY{p}{,} \PY{p}{[}\PY{o}{\PYZhy{}}\PY{l+m+mi}{1}\PY{p}{,} \PY{l+m+mi}{1}\PY{p}{]}\PY{p}{,} \PY{p}{[}\PY{l+m+mi}{0}\PY{p}{,} \PY{l+m+mi}{0}\PY{p}{]}\PY{p}{]}\PY{p}{)}
\PY{n}{word2ind\PYZus{}plot\PYZus{}test} \PY{o}{=} \PY{p}{\PYZob{}}\PY{l+s+s1}{\PYZsq{}}\PY{l+s+s1}{test1}\PY{l+s+s1}{\PYZsq{}}\PY{p}{:} \PY{l+m+mi}{0}\PY{p}{,} \PY{l+s+s1}{\PYZsq{}}\PY{l+s+s1}{test2}\PY{l+s+s1}{\PYZsq{}}\PY{p}{:} \PY{l+m+mi}{1}\PY{p}{,} \PY{l+s+s1}{\PYZsq{}}\PY{l+s+s1}{test3}\PY{l+s+s1}{\PYZsq{}}\PY{p}{:} \PY{l+m+mi}{2}\PY{p}{,} \PY{l+s+s1}{\PYZsq{}}\PY{l+s+s1}{test4}\PY{l+s+s1}{\PYZsq{}}\PY{p}{:} \PY{l+m+mi}{3}\PY{p}{,} \PY{l+s+s1}{\PYZsq{}}\PY{l+s+s1}{test5}\PY{l+s+s1}{\PYZsq{}}\PY{p}{:} \PY{l+m+mi}{4}\PY{p}{\PYZcb{}}
\PY{n}{words} \PY{o}{=} \PY{p}{[}\PY{l+s+s1}{\PYZsq{}}\PY{l+s+s1}{test1}\PY{l+s+s1}{\PYZsq{}}\PY{p}{,} \PY{l+s+s1}{\PYZsq{}}\PY{l+s+s1}{test2}\PY{l+s+s1}{\PYZsq{}}\PY{p}{,} \PY{l+s+s1}{\PYZsq{}}\PY{l+s+s1}{test3}\PY{l+s+s1}{\PYZsq{}}\PY{p}{,} \PY{l+s+s1}{\PYZsq{}}\PY{l+s+s1}{test4}\PY{l+s+s1}{\PYZsq{}}\PY{p}{,} \PY{l+s+s1}{\PYZsq{}}\PY{l+s+s1}{test5}\PY{l+s+s1}{\PYZsq{}}\PY{p}{]}
\PY{n}{plot\PYZus{}embeddings}\PY{p}{(}\PY{n}{M\PYZus{}reduced\PYZus{}plot\PYZus{}test}\PY{p}{,} \PY{n}{word2ind\PYZus{}plot\PYZus{}test}\PY{p}{,} \PY{n}{words}\PY{p}{)}

\PY{n+nb}{print} \PY{p}{(}\PY{l+s+s2}{\PYZdq{}}\PY{l+s+s2}{\PYZhy{}}\PY{l+s+s2}{\PYZdq{}} \PY{o}{*} \PY{l+m+mi}{80}\PY{p}{)}
\end{Verbatim}
\end{tcolorbox}

    \begin{Verbatim}[commandchars=\\\{\}]
--------------------------------------------------------------------------------
Outputted Plot:
    \end{Verbatim}

    \begin{center}
    \adjustimage{max size={0.9\linewidth}{0.9\paperheight}}{exploring_word_vectors_files/exploring_word_vectors_20_1.png}
    \end{center}
    { \hspace*{\fill} \\}
    
    \begin{Verbatim}[commandchars=\\\{\}]
--------------------------------------------------------------------------------
    \end{Verbatim}

    \subsubsection{Question 1.5: Co-Occurrence Plot Analysis {[}written{]}
(3
points)}\label{question-1.5-co-occurrence-plot-analysis-written-3-points}

Now we will put together all the parts you have written! We will compute
the co-occurrence matrix with fixed window of 4 (the default window
size), over the Large Movie Review corpus. Then we will use TruncatedSVD
to compute 2-dimensional embeddings of each word. TruncatedSVD returns
U*S, so we need to normalize the returned vectors, so that all the
vectors will appear around the unit circle (therefore closeness is
directional closeness). \textbf{Note}: The line of code below that does
the normalizing uses the NumPy concept of \emph{broadcasting}. If you
don't know about broadcasting, check out
\href{https://jakevdp.github.io/PythonDataScienceHandbook/02.05-computation-on-arrays-broadcasting.html}{Computation
on Arrays: Broadcasting by Jake VanderPlas}.

Run the below cell to produce the plot. It can take up to a few minutes
to run.

    \begin{tcolorbox}[breakable, size=fbox, boxrule=1pt, pad at break*=1mm,colback=cellbackground, colframe=cellborder]
\prompt{In}{incolor}{16}{\boxspacing}
\begin{Verbatim}[commandchars=\\\{\}]
\PY{c+c1}{\PYZsh{} \PYZhy{}\PYZhy{}\PYZhy{}\PYZhy{}\PYZhy{}\PYZhy{}\PYZhy{}\PYZhy{}\PYZhy{}\PYZhy{}\PYZhy{}\PYZhy{}\PYZhy{}\PYZhy{}\PYZhy{}\PYZhy{}\PYZhy{}\PYZhy{}\PYZhy{}\PYZhy{}\PYZhy{}\PYZhy{}\PYZhy{}\PYZhy{}\PYZhy{}\PYZhy{}\PYZhy{}\PYZhy{}\PYZhy{}}
\PY{c+c1}{\PYZsh{} Run This Cell to Produce Your Plot}
\PY{c+c1}{\PYZsh{} \PYZhy{}\PYZhy{}\PYZhy{}\PYZhy{}\PYZhy{}\PYZhy{}\PYZhy{}\PYZhy{}\PYZhy{}\PYZhy{}\PYZhy{}\PYZhy{}\PYZhy{}\PYZhy{}\PYZhy{}\PYZhy{}\PYZhy{}\PYZhy{}\PYZhy{}\PYZhy{}\PYZhy{}\PYZhy{}\PYZhy{}\PYZhy{}\PYZhy{}\PYZhy{}\PYZhy{}\PYZhy{}\PYZhy{}\PYZhy{}}
\PY{n}{imdb\PYZus{}corpus} \PY{o}{=} \PY{n}{read\PYZus{}corpus}\PY{p}{(}\PY{p}{)}
\PY{n}{M\PYZus{}co\PYZus{}occurrence}\PY{p}{,} \PY{n}{word2ind\PYZus{}co\PYZus{}occurrence} \PY{o}{=} \PY{n}{compute\PYZus{}co\PYZus{}occurrence\PYZus{}matrix}\PY{p}{(}\PY{n}{imdb\PYZus{}corpus}\PY{p}{)}
\PY{n}{M\PYZus{}reduced\PYZus{}co\PYZus{}occurrence} \PY{o}{=} \PY{n}{reduce\PYZus{}to\PYZus{}k\PYZus{}dim}\PY{p}{(}\PY{n}{M\PYZus{}co\PYZus{}occurrence}\PY{p}{,} \PY{n}{k}\PY{o}{=}\PY{l+m+mi}{2}\PY{p}{)}

\PY{c+c1}{\PYZsh{} Rescale (normalize) the rows to make them each of unit\PYZhy{}length}
\PY{n}{M\PYZus{}lengths} \PY{o}{=} \PY{n}{np}\PY{o}{.}\PY{n}{linalg}\PY{o}{.}\PY{n}{norm}\PY{p}{(}\PY{n}{M\PYZus{}reduced\PYZus{}co\PYZus{}occurrence}\PY{p}{,} \PY{n}{axis}\PY{o}{=}\PY{l+m+mi}{1}\PY{p}{)}
\PY{n}{M\PYZus{}normalized} \PY{o}{=} \PY{n}{M\PYZus{}reduced\PYZus{}co\PYZus{}occurrence} \PY{o}{/} \PY{n}{M\PYZus{}lengths}\PY{p}{[}\PY{p}{:}\PY{p}{,} \PY{n}{np}\PY{o}{.}\PY{n}{newaxis}\PY{p}{]} \PY{c+c1}{\PYZsh{} broadcasting}

\PY{n}{words} \PY{o}{=} \PY{p}{[}\PY{l+s+s1}{\PYZsq{}}\PY{l+s+s1}{movie}\PY{l+s+s1}{\PYZsq{}}\PY{p}{,} \PY{l+s+s1}{\PYZsq{}}\PY{l+s+s1}{book}\PY{l+s+s1}{\PYZsq{}}\PY{p}{,} \PY{l+s+s1}{\PYZsq{}}\PY{l+s+s1}{mysterious}\PY{l+s+s1}{\PYZsq{}}\PY{p}{,} \PY{l+s+s1}{\PYZsq{}}\PY{l+s+s1}{story}\PY{l+s+s1}{\PYZsq{}}\PY{p}{,} \PY{l+s+s1}{\PYZsq{}}\PY{l+s+s1}{fascinating}\PY{l+s+s1}{\PYZsq{}}\PY{p}{,} \PY{l+s+s1}{\PYZsq{}}\PY{l+s+s1}{good}\PY{l+s+s1}{\PYZsq{}}\PY{p}{,} \PY{l+s+s1}{\PYZsq{}}\PY{l+s+s1}{interesting}\PY{l+s+s1}{\PYZsq{}}\PY{p}{,} \PY{l+s+s1}{\PYZsq{}}\PY{l+s+s1}{large}\PY{l+s+s1}{\PYZsq{}}\PY{p}{,} \PY{l+s+s1}{\PYZsq{}}\PY{l+s+s1}{massive}\PY{l+s+s1}{\PYZsq{}}\PY{p}{,} \PY{l+s+s1}{\PYZsq{}}\PY{l+s+s1}{huge}\PY{l+s+s1}{\PYZsq{}}\PY{p}{]}

\PY{n}{plot\PYZus{}embeddings}\PY{p}{(}\PY{n}{M\PYZus{}normalized}\PY{p}{,} \PY{n}{word2ind\PYZus{}co\PYZus{}occurrence}\PY{p}{,} \PY{n}{words}\PY{p}{)}
\end{Verbatim}
\end{tcolorbox}

    \begin{Verbatim}[commandchars=\\\{\}]
Running Truncated SVD over 5880 words{\ldots}
Done.
    \end{Verbatim}

    \begin{center}
    \adjustimage{max size={0.9\linewidth}{0.9\paperheight}}{exploring_word_vectors_files/exploring_word_vectors_22_1.png}
    \end{center}
    { \hspace*{\fill} \\}
    
    \textbf{Verify that your figure matches ``question\_1.5.png'' in the
assignment zip. If not, use the figure in ``question\_1.5.png'' to
answer the next two questions.}

    \begin{enumerate}
\def\labelenumi{\alph{enumi}.}
\tightlist
\item
  Find at least two groups of words that cluster together in
  2-dimensional embedding space. Give an explanation for each cluster
  you observe.
\end{enumerate}

    \paragraph{Write your answer here.}\label{write-your-answer-here.}

fascinating、good、interesting。这三个词都是正面评价形容词，语义上都表示``好的、吸引人的、有趣的''。在语料中它们经常出现在相似的上下文中，例如``a
good movie''、``an interesting book''、``a fascinating
story''，因此它们的词向量在训练后会被拉近，在降维后的二维空间中也表现出紧密的聚类；

story、book、movie。这三个词都与叙事、媒体内容相关，是典型的共现词。例如``the
story of the book/movie''、``read a book''、``watch a
movie''，它们在语料中频繁共同出现，上下文高度重叠，因此嵌入向量相似，在二维投影中聚集在一起。

    \begin{enumerate}
\def\labelenumi{\alph{enumi}.}
\setcounter{enumi}{1}
\tightlist
\item
  What doesn't cluster together that you might think should have?
  Describe at least two examples.
\end{enumerate}

    \paragraph{Write your answer here.}\label{write-your-answer-here.}

massive和large、huge和massive应该聚在一起

    \subsection{Part 2: Prediction-Based Word Vectors (15
points)}\label{part-2-prediction-based-word-vectors-15-points}

As discussed in class, more recently prediction-based word vectors have
demonstrated better performance, such as word2vec and GloVe (which also
utilizes the benefit of counts). Here, we shall explore the embeddings
produced by GloVe. Please revisit the class notes and lecture slides for
more details on the word2vec and GloVe algorithms. If you're feeling
adventurous, challenge yourself and try reading
\href{https://nlp.stanford.edu/pubs/glove.pdf}{GloVe's original paper}.

Then run the following cells to load the GloVe vectors into memory.
\textbf{Note}: If this is your first time to run these cells,
i.e.~download the embedding model, it will take a couple minutes to run.
If you've run these cells before, rerunning them will load the model
without redownloading it, which will take about 1 to 2 minutes.

    \begin{tcolorbox}[breakable, size=fbox, boxrule=1pt, pad at break*=1mm,colback=cellbackground, colframe=cellborder]
\prompt{In}{incolor}{17}{\boxspacing}
\begin{Verbatim}[commandchars=\\\{\}]
\PY{k}{def}\PY{+w}{ }\PY{n+nf}{load\PYZus{}embedding\PYZus{}model}\PY{p}{(}\PY{p}{)}\PY{p}{:}
\PY{+w}{    }\PY{l+s+sd}{\PYZdq{}\PYZdq{}\PYZdq{} Load GloVe Vectors}
\PY{l+s+sd}{        Return:}
\PY{l+s+sd}{            wv\PYZus{}from\PYZus{}bin: All 400000 embeddings, each length 200}
\PY{l+s+sd}{    \PYZdq{}\PYZdq{}\PYZdq{}}
    \PY{k+kn}{import}\PY{+w}{ }\PY{n+nn}{gensim}\PY{n+nn}{.}\PY{n+nn}{downloader}\PY{+w}{ }\PY{k}{as}\PY{+w}{ }\PY{n+nn}{api}
    \PY{n}{wv\PYZus{}from\PYZus{}bin} \PY{o}{=} \PY{n}{api}\PY{o}{.}\PY{n}{load}\PY{p}{(}\PY{l+s+s2}{\PYZdq{}}\PY{l+s+s2}{glove\PYZhy{}wiki\PYZhy{}gigaword\PYZhy{}200}\PY{l+s+s2}{\PYZdq{}}\PY{p}{)}
    \PY{n+nb}{print}\PY{p}{(}\PY{l+s+s2}{\PYZdq{}}\PY{l+s+s2}{Loaded vocab size }\PY{l+s+si}{\PYZpc{}i}\PY{l+s+s2}{\PYZdq{}} \PY{o}{\PYZpc{}} \PY{n+nb}{len}\PY{p}{(}\PY{n+nb}{list}\PY{p}{(}\PY{n}{wv\PYZus{}from\PYZus{}bin}\PY{o}{.}\PY{n}{index\PYZus{}to\PYZus{}key}\PY{p}{)}\PY{p}{)}\PY{p}{)}
    \PY{k}{return} \PY{n}{wv\PYZus{}from\PYZus{}bin}
\PY{n}{wv\PYZus{}from\PYZus{}bin} \PY{o}{=} \PY{n}{load\PYZus{}embedding\PYZus{}model}\PY{p}{(}\PY{p}{)}
\end{Verbatim}
\end{tcolorbox}

    \begin{Verbatim}[commandchars=\\\{\}]
[==================================================] 100.0\% 252.1/252.1MB
downloaded
Loaded vocab size 400000
    \end{Verbatim}

    \paragraph{Note: If you are receiving a ``reset by peer'' error, rerun
the cell to restart the
download.}\label{note-if-you-are-receiving-a-reset-by-peer-error-rerun-the-cell-to-restart-the-download.}

    \subsubsection{Reducing dimensionality of Word
Embeddings}\label{reducing-dimensionality-of-word-embeddings}

Let's directly compare the GloVe embeddings to those of the
co-occurrence matrix. In order to avoid running out of memory, we will
work with a sample of 40000 GloVe vectors instead. Run the following
cells to:

\begin{enumerate}
\def\labelenumi{\arabic{enumi}.}
\tightlist
\item
  Put 40000 Glove vectors into a matrix M
\item
  Run \texttt{reduce\_to\_k\_dim} (your Truncated SVD function) to
  reduce the vectors from 200-dimensional to 2-dimensional.
\end{enumerate}

    \begin{tcolorbox}[breakable, size=fbox, boxrule=1pt, pad at break*=1mm,colback=cellbackground, colframe=cellborder]
\prompt{In}{incolor}{18}{\boxspacing}
\begin{Verbatim}[commandchars=\\\{\}]
\PY{k}{def}\PY{+w}{ }\PY{n+nf}{get\PYZus{}matrix\PYZus{}of\PYZus{}vectors}\PY{p}{(}\PY{n}{wv\PYZus{}from\PYZus{}bin}\PY{p}{,} \PY{n}{required\PYZus{}words}\PY{p}{)}\PY{p}{:}
\PY{+w}{    }\PY{l+s+sd}{\PYZdq{}\PYZdq{}\PYZdq{} Put the GloVe vectors into a matrix M.}
\PY{l+s+sd}{        Param:}
\PY{l+s+sd}{            wv\PYZus{}from\PYZus{}bin: KeyedVectors object; the 400000 GloVe vectors loaded from file}
\PY{l+s+sd}{        Return:}
\PY{l+s+sd}{            M: numpy matrix shape (num words, 200) containing the vectors}
\PY{l+s+sd}{            word2ind: dictionary mapping each word to its row number in M}
\PY{l+s+sd}{    \PYZdq{}\PYZdq{}\PYZdq{}}
    \PY{k+kn}{import}\PY{+w}{ }\PY{n+nn}{random}
    \PY{n}{words} \PY{o}{=} \PY{n+nb}{list}\PY{p}{(}\PY{n}{wv\PYZus{}from\PYZus{}bin}\PY{o}{.}\PY{n}{index\PYZus{}to\PYZus{}key}\PY{p}{)}
    \PY{n+nb}{print}\PY{p}{(}\PY{l+s+s2}{\PYZdq{}}\PY{l+s+s2}{Shuffling words ...}\PY{l+s+s2}{\PYZdq{}}\PY{p}{)}
    \PY{n}{random}\PY{o}{.}\PY{n}{seed}\PY{p}{(}\PY{l+m+mi}{225}\PY{p}{)}
    \PY{n}{random}\PY{o}{.}\PY{n}{shuffle}\PY{p}{(}\PY{n}{words}\PY{p}{)}
    \PY{n+nb}{print}\PY{p}{(}\PY{l+s+s2}{\PYZdq{}}\PY{l+s+s2}{Putting }\PY{l+s+si}{\PYZpc{}i}\PY{l+s+s2}{ words into word2ind and matrix M...}\PY{l+s+s2}{\PYZdq{}} \PY{o}{\PYZpc{}} \PY{n+nb}{len}\PY{p}{(}\PY{n}{words}\PY{p}{)}\PY{p}{)}
    \PY{n}{word2ind} \PY{o}{=} \PY{p}{\PYZob{}}\PY{p}{\PYZcb{}}
    \PY{n}{M} \PY{o}{=} \PY{p}{[}\PY{p}{]}
    \PY{n}{curInd} \PY{o}{=} \PY{l+m+mi}{0}
    \PY{k}{for} \PY{n}{w} \PY{o+ow}{in} \PY{n}{words}\PY{p}{:}
        \PY{k}{try}\PY{p}{:}
            \PY{n}{M}\PY{o}{.}\PY{n}{append}\PY{p}{(}\PY{n}{wv\PYZus{}from\PYZus{}bin}\PY{o}{.}\PY{n}{get\PYZus{}vector}\PY{p}{(}\PY{n}{w}\PY{p}{)}\PY{p}{)}
            \PY{n}{word2ind}\PY{p}{[}\PY{n}{w}\PY{p}{]} \PY{o}{=} \PY{n}{curInd}
            \PY{n}{curInd} \PY{o}{+}\PY{o}{=} \PY{l+m+mi}{1}
        \PY{k}{except} \PY{n+ne}{KeyError}\PY{p}{:}
            \PY{k}{continue}
    \PY{k}{for} \PY{n}{w} \PY{o+ow}{in} \PY{n}{required\PYZus{}words}\PY{p}{:}
        \PY{k}{if} \PY{n}{w} \PY{o+ow}{in} \PY{n}{words}\PY{p}{:}
            \PY{k}{continue}
        \PY{k}{try}\PY{p}{:}
            \PY{n}{M}\PY{o}{.}\PY{n}{append}\PY{p}{(}\PY{n}{wv\PYZus{}from\PYZus{}bin}\PY{o}{.}\PY{n}{get\PYZus{}vector}\PY{p}{(}\PY{n}{w}\PY{p}{)}\PY{p}{)}
            \PY{n}{word2ind}\PY{p}{[}\PY{n}{w}\PY{p}{]} \PY{o}{=} \PY{n}{curInd}
            \PY{n}{curInd} \PY{o}{+}\PY{o}{=} \PY{l+m+mi}{1}
        \PY{k}{except} \PY{n+ne}{KeyError}\PY{p}{:}
            \PY{k}{continue}
    \PY{n}{M} \PY{o}{=} \PY{n}{np}\PY{o}{.}\PY{n}{stack}\PY{p}{(}\PY{n}{M}\PY{p}{)}
    \PY{n+nb}{print}\PY{p}{(}\PY{l+s+s2}{\PYZdq{}}\PY{l+s+s2}{Done.}\PY{l+s+s2}{\PYZdq{}}\PY{p}{)}
    \PY{k}{return} \PY{n}{M}\PY{p}{,} \PY{n}{word2ind}
\end{Verbatim}
\end{tcolorbox}

    \begin{tcolorbox}[breakable, size=fbox, boxrule=1pt, pad at break*=1mm,colback=cellbackground, colframe=cellborder]
\prompt{In}{incolor}{19}{\boxspacing}
\begin{Verbatim}[commandchars=\\\{\}]
\PY{c+c1}{\PYZsh{} \PYZhy{}\PYZhy{}\PYZhy{}\PYZhy{}\PYZhy{}\PYZhy{}\PYZhy{}\PYZhy{}\PYZhy{}\PYZhy{}\PYZhy{}\PYZhy{}\PYZhy{}\PYZhy{}\PYZhy{}\PYZhy{}\PYZhy{}\PYZhy{}\PYZhy{}\PYZhy{}\PYZhy{}\PYZhy{}\PYZhy{}\PYZhy{}\PYZhy{}\PYZhy{}\PYZhy{}\PYZhy{}\PYZhy{}\PYZhy{}\PYZhy{}\PYZhy{}\PYZhy{}\PYZhy{}\PYZhy{}\PYZhy{}\PYZhy{}\PYZhy{}\PYZhy{}\PYZhy{}\PYZhy{}\PYZhy{}\PYZhy{}\PYZhy{}\PYZhy{}\PYZhy{}\PYZhy{}\PYZhy{}\PYZhy{}\PYZhy{}\PYZhy{}\PYZhy{}\PYZhy{}\PYZhy{}\PYZhy{}\PYZhy{}\PYZhy{}\PYZhy{}\PYZhy{}\PYZhy{}\PYZhy{}\PYZhy{}\PYZhy{}\PYZhy{}\PYZhy{}}
\PY{c+c1}{\PYZsh{} Run Cell to Reduce 200\PYZhy{}Dimensional Word Embeddings to k Dimensions}
\PY{c+c1}{\PYZsh{} Note: This should be quick to run}
\PY{c+c1}{\PYZsh{} \PYZhy{}\PYZhy{}\PYZhy{}\PYZhy{}\PYZhy{}\PYZhy{}\PYZhy{}\PYZhy{}\PYZhy{}\PYZhy{}\PYZhy{}\PYZhy{}\PYZhy{}\PYZhy{}\PYZhy{}\PYZhy{}\PYZhy{}\PYZhy{}\PYZhy{}\PYZhy{}\PYZhy{}\PYZhy{}\PYZhy{}\PYZhy{}\PYZhy{}\PYZhy{}\PYZhy{}\PYZhy{}\PYZhy{}\PYZhy{}\PYZhy{}\PYZhy{}\PYZhy{}\PYZhy{}\PYZhy{}\PYZhy{}\PYZhy{}\PYZhy{}\PYZhy{}\PYZhy{}\PYZhy{}\PYZhy{}\PYZhy{}\PYZhy{}\PYZhy{}\PYZhy{}\PYZhy{}\PYZhy{}\PYZhy{}\PYZhy{}\PYZhy{}\PYZhy{}\PYZhy{}\PYZhy{}\PYZhy{}\PYZhy{}\PYZhy{}\PYZhy{}\PYZhy{}\PYZhy{}\PYZhy{}\PYZhy{}\PYZhy{}\PYZhy{}\PYZhy{}}
\PY{n}{M}\PY{p}{,} \PY{n}{word2ind} \PY{o}{=} \PY{n}{get\PYZus{}matrix\PYZus{}of\PYZus{}vectors}\PY{p}{(}\PY{n}{wv\PYZus{}from\PYZus{}bin}\PY{p}{,} \PY{n}{words}\PY{p}{)}
\PY{n}{M\PYZus{}reduced} \PY{o}{=} \PY{n}{reduce\PYZus{}to\PYZus{}k\PYZus{}dim}\PY{p}{(}\PY{n}{M}\PY{p}{,} \PY{n}{k}\PY{o}{=}\PY{l+m+mi}{2}\PY{p}{)}

\PY{c+c1}{\PYZsh{} Rescale (normalize) the rows to make them each of unit\PYZhy{}length}
\PY{n}{M\PYZus{}lengths} \PY{o}{=} \PY{n}{np}\PY{o}{.}\PY{n}{linalg}\PY{o}{.}\PY{n}{norm}\PY{p}{(}\PY{n}{M\PYZus{}reduced}\PY{p}{,} \PY{n}{axis}\PY{o}{=}\PY{l+m+mi}{1}\PY{p}{)}
\PY{n}{M\PYZus{}reduced\PYZus{}normalized} \PY{o}{=} \PY{n}{M\PYZus{}reduced} \PY{o}{/} \PY{n}{M\PYZus{}lengths}\PY{p}{[}\PY{p}{:}\PY{p}{,} \PY{n}{np}\PY{o}{.}\PY{n}{newaxis}\PY{p}{]} \PY{c+c1}{\PYZsh{} broadcasting}
\end{Verbatim}
\end{tcolorbox}

    \begin{Verbatim}[commandchars=\\\{\}]
Shuffling words {\ldots}
Putting 400000 words into word2ind and matrix M{\ldots}
Done.
Running Truncated SVD over 400000 words{\ldots}
Done.
    \end{Verbatim}

    \textbf{Note: If you are receiving out of memory issues on your local
machine, try closing other applications to free more memory on your
device. You may want to try restarting your machine so that you can free
up extra memory. Then immediately run the jupyter notebook and see if
you can load the word vectors properly. If you still have problems with
loading the embeddings onto your local machine after this, please go to
office hours or contact course staff.}

    \subsubsection{Question 2.1: GloVe Plot Analysis {[}written{]} (3
points)}\label{question-2.1-glove-plot-analysis-written-3-points}

Run the cell below to plot the 2D GloVe embeddings for
\texttt{{[}\textquotesingle{}movie\textquotesingle{},\ \textquotesingle{}book\textquotesingle{},\ \textquotesingle{}mysterious\textquotesingle{},\ \textquotesingle{}story\textquotesingle{},\ \textquotesingle{}fascinating\textquotesingle{},\ \textquotesingle{}good\textquotesingle{},\ \textquotesingle{}interesting\textquotesingle{},\ \textquotesingle{}large\textquotesingle{},\ \textquotesingle{}massive\textquotesingle{},\ \textquotesingle{}huge\textquotesingle{}{]}}.

    \begin{tcolorbox}[breakable, size=fbox, boxrule=1pt, pad at break*=1mm,colback=cellbackground, colframe=cellborder]
\prompt{In}{incolor}{20}{\boxspacing}
\begin{Verbatim}[commandchars=\\\{\}]
\PY{n}{words} \PY{o}{=} \PY{p}{[}\PY{l+s+s1}{\PYZsq{}}\PY{l+s+s1}{movie}\PY{l+s+s1}{\PYZsq{}}\PY{p}{,} \PY{l+s+s1}{\PYZsq{}}\PY{l+s+s1}{book}\PY{l+s+s1}{\PYZsq{}}\PY{p}{,} \PY{l+s+s1}{\PYZsq{}}\PY{l+s+s1}{mysterious}\PY{l+s+s1}{\PYZsq{}}\PY{p}{,} \PY{l+s+s1}{\PYZsq{}}\PY{l+s+s1}{story}\PY{l+s+s1}{\PYZsq{}}\PY{p}{,} \PY{l+s+s1}{\PYZsq{}}\PY{l+s+s1}{fascinating}\PY{l+s+s1}{\PYZsq{}}\PY{p}{,} \PY{l+s+s1}{\PYZsq{}}\PY{l+s+s1}{good}\PY{l+s+s1}{\PYZsq{}}\PY{p}{,} \PY{l+s+s1}{\PYZsq{}}\PY{l+s+s1}{interesting}\PY{l+s+s1}{\PYZsq{}}\PY{p}{,} \PY{l+s+s1}{\PYZsq{}}\PY{l+s+s1}{large}\PY{l+s+s1}{\PYZsq{}}\PY{p}{,} \PY{l+s+s1}{\PYZsq{}}\PY{l+s+s1}{massive}\PY{l+s+s1}{\PYZsq{}}\PY{p}{,} \PY{l+s+s1}{\PYZsq{}}\PY{l+s+s1}{huge}\PY{l+s+s1}{\PYZsq{}}\PY{p}{]}

\PY{n}{plot\PYZus{}embeddings}\PY{p}{(}\PY{n}{M\PYZus{}reduced\PYZus{}normalized}\PY{p}{,} \PY{n}{word2ind}\PY{p}{,} \PY{n}{words}\PY{p}{)}
\end{Verbatim}
\end{tcolorbox}

    \begin{center}
    \adjustimage{max size={0.9\linewidth}{0.9\paperheight}}{exploring_word_vectors_files/exploring_word_vectors_36_0.png}
    \end{center}
    { \hspace*{\fill} \\}
    
    \textbf{Verify that your figure matches ``question\_2.1.png'' in the
assignment zip. If not, use the figure in ``question\_2.1.png'' (and the
figure in ``question\_1.5.png'', if applicable) to answer the next two
questions.}

    \begin{enumerate}
\def\labelenumi{\alph{enumi}.}
\tightlist
\item
  What is one way the plot is different from the one generated earlier
  from the co-occurrence matrix? What is one way it's similar?
\end{enumerate}

    \paragraph{Write your answer here.}\label{write-your-answer-here.}

不同点：在共现矩阵+SVD图中，词语是按词性或语义大类分群的。例如，fascinating、good、interesting
三个评价性形容词紧紧聚在一起；story、book、movie
三个名词聚在一起；大小形容词（massive, large,
huge）虽然分散，但大致散落在不同角落。

而在GloVe图中，这种分群被打破了。good 脱离了形容词群，跑到了右上角，和
book、story、movie 挤在了一起；interesting、mysterious、huge
这三个词聚在了一堆；而 fascinating
被完全孤立在左下角。这说明GloVe捕捉到的语义相似性，与直接分解共现矩阵捕捉到的相似性有了本质区别。

相似点：无论在哪张图中，book、story、movie
始终紧密地聚在一起，形成整个图中最稳固的核心簇。同时，massive 和 large
在两张图中都相对靠近（在GloVe图中两者距离更近）。这说明尽管方法不同，它们都成功捕捉到了``叙事媒介名词''和``巨大''这两个强语义关联。

    \begin{enumerate}
\def\labelenumi{\alph{enumi}.}
\setcounter{enumi}{1}
\tightlist
\item
  Why might the GloVe plot (question\_2.1.png) differ from the plot
  generated earlier from the co-occurrence matrix (question\_1.5.png)?
\end{enumerate}

    \paragraph{Write your answer here.}\label{write-your-answer-here.}

目标函数不同（最核心的原因）：

\begin{verbatim}
共现矩阵+SVD (LSA)：目标是矩阵重构。它试图用低维向量尽可能还原原始共现矩阵的数值大小（最小化Frobenius范数）。因此，它保留的是全局的方差分布，这导致高频共现的词（即使语义类别不同）容易聚在一起。

GloVe：目标是拟合共现概率的对数。GloVe的核心创新是引入了共现概率的比值（Ratio of Co-occurrence Probabilities）。例如，“冰”和“蒸汽”的共现概率之比，能更好地帮模型区分“水”（固体）和“水”（气体）。这种机制使得GloVe对区分性特征（即词与词之间的细微语义差异）更敏感。因此，good 可能因为经常直接修饰 book（如 "good book"）而在GloVe中与 book 靠得更近；而SVD则因为 good 在全局上广泛作为评价词，将其与 interesting 归为一类。
\end{verbatim}

上下文窗口与加权机制不同：

\begin{verbatim}
LSA通常使用固定窗口或整篇文档的原始共现计数（或TF-IDF），所有共现一视同仁。

GloVe在构建共现矩阵时，引入了距离衰减加权（通常为 1/d）。距离越近的词对，权重越大。这改变了词向量的相似度计算，使得近义词（如 massive 和 large，它们经常在相近位置互相替代或修饰同一名词）更容易被拉近，而 huge 可能因为搭配习惯不同（如 "huge success"）被推向了别的簇。
\end{verbatim}

降维可视化步骤的二次失真：

\begin{verbatim}
共现矩阵图通常是直接对共现矩阵做SVD降到2维，或者对高维SVD结果做截断。

GloVe图通常是先在较高维度（如50维、100维）训练出词向量，然后使用 t-SNE、UMAP 或 PCA 将其再次降维到2维进行可视化。t-SNE等非线性降维极度偏好保留局部邻域，这会放大局部聚集（例如把 interesting、mysterious、huge 强行拉在一起），但会破坏全局的语义大类（例如把 fascinating 远远甩开）。这个二次降维的过程极大地扭曲了高维空间中的真实距离。
\end{verbatim}

优化算法与随机性：

\begin{verbatim}
SVD是确定性的线性代数分解（即使有随机化SVD，也是矩阵运算）。

GloVe是使用随机梯度下降（AdaGrad）训练的非凸优化模型。权重的随机初始化、学习率、迭代次数以及 random_state 的不同，都会导致收敛到不同的局部最优解。这种随机性导致其空间方向与SVD空间没有直接的可比性。
\end{verbatim}

    \subsubsection{Cosine Similarity}\label{cosine-similarity}

Now that we have word vectors, we need a way to quantify the similarity
between individual words, according to these vectors. One such metric is
cosine-similarity. We will be using this to find words that are
``close'' and ``far'' from one another.

We can think of n-dimensional vectors as points in n-dimensional space.
If we take this perspective
\href{http://mathworld.wolfram.com/L1-Norm.html}{L1} and
\href{http://mathworld.wolfram.com/L2-Norm.html}{L2} Distances help
quantify the amount of space ``we must travel'' to get between these two
points. Another approach is to examine the angle between two vectors.
From trigonometry we know that:

Instead of computing the actual angle, we can leave the similarity in
terms of \(similarity = cos(\Theta)\). Formally the
\href{https://en.wikipedia.org/wiki/Cosine_similarity}{Cosine
Similarity} \(s\) between two vectors \(p\) and \(q\) is defined as:

\[s = \frac{p \cdot q}{||p|| ||q||}, \textrm{ where } s \in [-1, 1] \]

    \subsubsection{Question 2.2: Words with Multiple Meanings (1.5 points)
{[}code +
written{]}}\label{question-2.2-words-with-multiple-meanings-1.5-points-code-written}

Polysemes and homonyms are words that have more than one meaning (see
this \href{https://en.wikipedia.org/wiki/Polysemy}{wiki page} to learn
more about the difference between polysemes and homonyms ). Find a word
with \emph{at least two different meanings} such that the top-10 most
similar words (according to cosine similarity) contain related words
from \emph{both} meanings. For example, ``leaves'' has both ``go\_away''
and ``a\_structure\_of\_a\_plant'' meaning in the top 10, and ``scoop''
has both ``handed\_waffle\_cone'' and ``lowdown''. You will probably
need to try several polysemous or homonymic words before you find one.

Please state the word you discover and the multiple meanings that occur
in the top 10. Why do you think many of the polysemous or homonymic
words you tried didn't work (i.e.~the top-10 most similar words only
contain \textbf{one} of the meanings of the words)?

\textbf{Note}: You should use the
\texttt{wv\_from\_bin.most\_similar(word)} function to get the top 10
most similar words. This function ranks all other words in the
vocabulary with respect to their cosine similarity to the given word.
For further assistance, please check the
\textbf{\href{https://radimrehurek.com/gensim/models/keyedvectors.html\#gensim.models.keyedvectors.FastTextKeyedVectors.most_similar}{GenSim
documentation}}.

    \begin{tcolorbox}[breakable, size=fbox, boxrule=1pt, pad at break*=1mm,colback=cellbackground, colframe=cellborder]
\prompt{In}{incolor}{21}{\boxspacing}
\begin{Verbatim}[commandchars=\\\{\}]
\PY{c+c1}{\PYZsh{} \PYZhy{}\PYZhy{}\PYZhy{}\PYZhy{}\PYZhy{}\PYZhy{}\PYZhy{}\PYZhy{}\PYZhy{}\PYZhy{}\PYZhy{}\PYZhy{}\PYZhy{}\PYZhy{}\PYZhy{}\PYZhy{}\PYZhy{}\PYZhy{}}
\PY{c+c1}{\PYZsh{} Write your implementation here.}

\PY{c+c1}{\PYZsh{} 使用 wv\PYZus{}from\PYZus{}bin.most\PYZus{}similar 获取前 10 个最相似词}
\PY{n}{wv\PYZus{}from\PYZus{}bin}\PY{o}{.}\PY{n}{most\PYZus{}similar}\PY{p}{(}\PY{l+s+s1}{\PYZsq{}}\PY{l+s+s1}{plant}\PY{l+s+s1}{\PYZsq{}}\PY{p}{)}
\PY{c+c1}{\PYZsh{} \PYZhy{}\PYZhy{}\PYZhy{}\PYZhy{}\PYZhy{}\PYZhy{}\PYZhy{}\PYZhy{}\PYZhy{}\PYZhy{}\PYZhy{}\PYZhy{}\PYZhy{}\PYZhy{}\PYZhy{}\PYZhy{}\PYZhy{}\PYZhy{}}
\end{Verbatim}
\end{tcolorbox}

            \begin{tcolorbox}[breakable, size=fbox, boxrule=.5pt, pad at break*=1mm, opacityfill=0]
\prompt{Out}{outcolor}{21}{\boxspacing}
\begin{Verbatim}[commandchars=\\\{\}]
[('plants', 0.8802077174186707),
 ('factory', 0.6054627299308777),
 ('facility', 0.556488573551178),
 ('produce', 0.5439175367355347),
 ('production', 0.5205835103988647),
 ('chemical', 0.5185851454734802),
 ('producing', 0.5161451101303101),
 ('waste', 0.5101147890090942),
 ('factories', 0.506610631942749),
 ('flowering', 0.5009768009185791)]
\end{Verbatim}
\end{tcolorbox}
        
    \paragraph{Write your answer here.}\label{write-your-answer-here.}

在 Word2Vec、GloVe
等传统词嵌入模型中，每个词只被分配一个单一的向量表示，无论它有多少种含义。模型在训练时，会将一个词在所有上下文中的共现信息全部``压缩''到一个向量中。

当多义词的不同含义在语料中使用频率极不均衡时，问题就会凸显：

\begin{verbatim}
主导含义压倒次要含义：如果某个含义在训练语料中出现得远多于其他含义，该词的向量就会主要反映主导含义。次要含义的共现信号会被淹没，导致 most_similar 返回的邻居词几乎全部来自主导含义。

语义空间中的“拉扯”效应：一个多义词的向量实际上是其所有含义的“加权平均”。如果不同含义的上下文差异很大（如 bank 的“河岸”和“银行”），这个平均向量可能落在两个语义簇的中间地带，最终与任何一方的相似度都不够高，导致 most_similar 返回的结果模糊或偏向一侧。

语料覆盖不足：如果训练语料中某个含义的样本很少，模型根本无法为其学习到有意义的表示。即使该含义在现实世界中很常见，只要语料中没有足够的相关上下文，词向量中就不会包含该含义的信息。
\end{verbatim}

    \subsubsection{Question 2.3: Synonyms \& Antonyms (2 points) {[}code +
written{]}}\label{question-2.3-synonyms-antonyms-2-points-code-written}

When considering Cosine Similarity, it's often more convenient to think
of Cosine Distance, which is simply 1 - Cosine Similarity.

Find three words \((w_1,w_2,w_3)\) where \(w_1\) and \(w_2\) are
synonyms and \(w_1\) and \(w_3\) are antonyms, but Cosine Distance
\((w_1,w_3) <\) Cosine Distance \((w_1,w_2)\).

As an example, \(w_1\)=``happy'' is closer to \(w_3\)=``sad'' than to
\(w_2\)=``cheerful''. Please find a different example that satisfies the
above. Once you have found your example, please give a possible
explanation for why this counter-intuitive result may have happened.

You should use the the \texttt{wv\_from\_bin.distance(w1,\ w2)} function
here in order to compute the cosine distance between two words. Please
see the
\textbf{\href{https://radimrehurek.com/gensim/models/keyedvectors.html\#gensim.models.keyedvectors.FastTextKeyedVectors.distance}{GenSim
documentation}} for further assistance.

    \begin{tcolorbox}[breakable, size=fbox, boxrule=1pt, pad at break*=1mm,colback=cellbackground, colframe=cellborder]
\prompt{In}{incolor}{22}{\boxspacing}
\begin{Verbatim}[commandchars=\\\{\}]
\PY{c+c1}{\PYZsh{} \PYZhy{}\PYZhy{}\PYZhy{}\PYZhy{}\PYZhy{}\PYZhy{}\PYZhy{}\PYZhy{}\PYZhy{}\PYZhy{}\PYZhy{}\PYZhy{}\PYZhy{}\PYZhy{}\PYZhy{}\PYZhy{}\PYZhy{}\PYZhy{}}
\PY{c+c1}{\PYZsh{} Write your implementation here.}

\PY{c+c1}{\PYZsh{} 使用 wv\PYZus{}from\PYZus{}bin.distance() 计算余弦距离}
\PY{n}{w1} \PY{o}{=} \PY{l+s+s2}{\PYZdq{}}\PY{l+s+s2}{happy}\PY{l+s+s2}{\PYZdq{}}
\PY{n}{w2} \PY{o}{=} \PY{l+s+s2}{\PYZdq{}}\PY{l+s+s2}{joyful}\PY{l+s+s2}{\PYZdq{}}   \PY{c+c1}{\PYZsh{} 同义词}
\PY{n}{w3} \PY{o}{=} \PY{l+s+s2}{\PYZdq{}}\PY{l+s+s2}{sad}\PY{l+s+s2}{\PYZdq{}}      \PY{c+c1}{\PYZsh{} 反义词}

\PY{n}{w1\PYZus{}w2\PYZus{}dist} \PY{o}{=} \PY{n}{wv\PYZus{}from\PYZus{}bin}\PY{o}{.}\PY{n}{distance}\PY{p}{(}\PY{n}{w1}\PY{p}{,} \PY{n}{w2}\PY{p}{)}
\PY{n}{w1\PYZus{}w3\PYZus{}dist} \PY{o}{=} \PY{n}{wv\PYZus{}from\PYZus{}bin}\PY{o}{.}\PY{n}{distance}\PY{p}{(}\PY{n}{w1}\PY{p}{,} \PY{n}{w3}\PY{p}{)}

\PY{n+nb}{print}\PY{p}{(}\PY{l+s+sa}{f}\PY{l+s+s2}{\PYZdq{}}\PY{l+s+s2}{Synonyms }\PY{l+s+s2}{\PYZsq{}}\PY{l+s+si}{\PYZob{}}\PY{n}{w1}\PY{l+s+si}{\PYZcb{}}\PY{l+s+s2}{\PYZsq{}}\PY{l+s+s2}{ and }\PY{l+s+s2}{\PYZsq{}}\PY{l+s+si}{\PYZob{}}\PY{n}{w2}\PY{l+s+si}{\PYZcb{}}\PY{l+s+s2}{\PYZsq{}}\PY{l+s+s2}{ have cosine distance: }\PY{l+s+si}{\PYZob{}}\PY{n}{w1\PYZus{}w2\PYZus{}dist}\PY{l+s+si}{\PYZcb{}}\PY{l+s+s2}{\PYZdq{}}\PY{p}{)}
\PY{n+nb}{print}\PY{p}{(}\PY{l+s+sa}{f}\PY{l+s+s2}{\PYZdq{}}\PY{l+s+s2}{Antonyms }\PY{l+s+s2}{\PYZsq{}}\PY{l+s+si}{\PYZob{}}\PY{n}{w1}\PY{l+s+si}{\PYZcb{}}\PY{l+s+s2}{\PYZsq{}}\PY{l+s+s2}{ and }\PY{l+s+s2}{\PYZsq{}}\PY{l+s+si}{\PYZob{}}\PY{n}{w3}\PY{l+s+si}{\PYZcb{}}\PY{l+s+s2}{\PYZsq{}}\PY{l+s+s2}{ have cosine distance: }\PY{l+s+si}{\PYZob{}}\PY{n}{w1\PYZus{}w3\PYZus{}dist}\PY{l+s+si}{\PYZcb{}}\PY{l+s+s2}{\PYZdq{}}\PY{p}{)}
\PY{c+c1}{\PYZsh{} \PYZhy{}\PYZhy{}\PYZhy{}\PYZhy{}\PYZhy{}\PYZhy{}\PYZhy{}\PYZhy{}\PYZhy{}\PYZhy{}\PYZhy{}\PYZhy{}\PYZhy{}\PYZhy{}\PYZhy{}\PYZhy{}\PYZhy{}\PYZhy{}}
\end{Verbatim}
\end{tcolorbox}

    \begin{Verbatim}[commandchars=\\\{\}]
Synonyms 'happy' and 'joyful' have cosine distance: 0.48148053884506226
Antonyms 'happy' and 'sad' have cosine distance: 0.4040136933326721
    \end{Verbatim}

    \paragraph{Write your answer here.}\label{write-your-answer-here.}

\begin{enumerate}
\def\labelenumi{\arabic{enumi}.}
\tightlist
\item
  分布假设：意义相近的词出现在相似的上下文中
\end{enumerate}

分布语义学的基本假设是：词义相似的词，其上下文也相似。然而，反义词恰恰经常出现在几乎相同的上下文中：

\begin{verbatim}
“The weather is hot / cold” —— hot 和 cold 出现的位置完全相同

“good movie” / “bad movie” —— good 和 bad 修饰相同的名词

“I feel happy today” / “I feel sad today” —— happy 和 sad 的句法环境几乎一致
\end{verbatim}

从共现统计的角度看，happy 和 sad 拥有高度重叠的上下文分布（都出现在``I
feel \textbf{\emph{''、``so }}''、``very
\_\_\_''等模式中），因此模型学到的向量会非常接近。 2.
模型无法捕捉``意义的方向''

Word2Vec 和 GloVe 只建模共现概率，而不建模语义极性。模型知道 happy 和
sad
出现在相似的语境中，但它无法理解这种相似性背后是``情感极性相反''还是``含义相同''。正如
StackOverflow 上的一则讨论所指出的：

\begin{verbatim}
“一个词和它的反义词经常会出现在彼此附近，以作文学对比或其他强调”
\end{verbatim}

\begin{enumerate}
\def\labelenumi{\arabic{enumi}.}
\setcounter{enumi}{2}
\tightlist
\item
  同义词反而可能因用法差异而距离更远
\end{enumerate}

happy 和 joyful
虽然语义相近，但在语料中的使用场景和搭配习惯可能存在差异：

\begin{verbatim}
happy 是一个高频、通用的词，搭配范围极广（happy birthday, happy hour, happy ending...）

joyful 相对低频，更多出现在文学性或正式语境中（joyful noise, joyful occasion）
\end{verbatim}

由于 joyful 的使用范围更窄，其上下文分布与 happy 的重叠程度反而可能低于
sad 与 happy 的重叠程度。happy 和 sad
作为情感光谱的两端，在``表达情绪''这一核心语境中的用法高度对称，而
joyful 可能还带有``喜悦、欢庆''的额外语义色彩，导致向量偏移。 4.
这是静态词嵌入的已知局限

多项研究指出，静态词嵌入在区分同义词和反义词方面存在系统性困难：

\begin{verbatim}
GloVe 的同义词平均余弦相似度约为 0.30，而反义词的平均余弦相似度反而更高，达到 0.45

对于名词而言，反义词比同义词更接近的情况甚至超过一半

有研究者直言：“高余弦相似度与同义关系相关的这一基础假设，已被证明是错误的”
\end{verbatim}

根本原因在于：静态词嵌入为每个词分配唯一一个向量，无论该词在不同上下文中含义如何变化。模型只能在``聚合所有用法''的层面上寻找相似性，而反义词由于共享大量上下文模式，反而在这种聚合视角下显得更``相似''。

    \subsubsection{Question 2.4: Analogies with Word Vectors {[}written{]}
(1.5
points)}\label{question-2.4-analogies-with-word-vectors-written-1.5-points}

Word vectors have been shown to \emph{sometimes} exhibit the ability to
solve analogies.

As an example, for the analogy ``man : grandfather :: woman : x'' (read:
man is to grandfather as woman is to x), what is x?

In the cell below, we show you how to use word vectors to find x using
the \texttt{most\_similar} function from the
\textbf{\href{https://radimrehurek.com/gensim/models/keyedvectors.html\#gensim.models.keyedvectors.KeyedVectors.most_similar}{GenSim
documentation}}. The function finds words that are most similar to the
words in the \texttt{positive} list and most dissimilar from the words
in the \texttt{negative} list (while omitting the input words, which are
often the most similar; see
\href{https://www.aclweb.org/anthology/N18-2039.pdf}{this paper}). The
answer to the analogy will have the highest cosine similarity (largest
returned numerical value).

    \begin{tcolorbox}[breakable, size=fbox, boxrule=1pt, pad at break*=1mm,colback=cellbackground, colframe=cellborder]
\prompt{In}{incolor}{23}{\boxspacing}
\begin{Verbatim}[commandchars=\\\{\}]
\PY{c+c1}{\PYZsh{} Run this cell to answer the analogy \PYZhy{}\PYZhy{} man : grandfather :: woman : x}
\PY{n}{pprint}\PY{o}{.}\PY{n}{pprint}\PY{p}{(}\PY{n}{wv\PYZus{}from\PYZus{}bin}\PY{o}{.}\PY{n}{most\PYZus{}similar}\PY{p}{(}\PY{n}{positive}\PY{o}{=}\PY{p}{[}\PY{l+s+s1}{\PYZsq{}}\PY{l+s+s1}{woman}\PY{l+s+s1}{\PYZsq{}}\PY{p}{,} \PY{l+s+s1}{\PYZsq{}}\PY{l+s+s1}{grandfather}\PY{l+s+s1}{\PYZsq{}}\PY{p}{]}\PY{p}{,} \PY{n}{negative}\PY{o}{=}\PY{p}{[}\PY{l+s+s1}{\PYZsq{}}\PY{l+s+s1}{man}\PY{l+s+s1}{\PYZsq{}}\PY{p}{]}\PY{p}{)}\PY{p}{)}
\end{Verbatim}
\end{tcolorbox}

    \begin{Verbatim}[commandchars=\\\{\}]
[('grandmother', 0.7608445286750793),
 ('granddaughter', 0.7200808525085449),
 ('daughter', 0.7168302536010742),
 ('mother', 0.7151536345481873),
 ('niece', 0.7005682587623596),
 ('father', 0.6659888029098511),
 ('aunt', 0.6623409390449524),
 ('grandson', 0.6618767976760864),
 ('grandparents', 0.644661009311676),
 ('wife', 0.6445354223251343)]
    \end{Verbatim}

    Let \(m\), \(g\), \(w\), and \(x\) denote the word vectors for
\texttt{man}, \texttt{grandfather}, \texttt{woman}, and the answer,
respectively. Using \textbf{only} vectors \(m\), \(g\), \(w\), and the
vector arithmetic operators \(+\) and \(-\) in your answer, what is the
expression in which we are maximizing cosine similarity with \(x\)?

Hint: Recall that word vectors are simply multi-dimensional vectors that
represent a word. It might help to draw out a 2D example using arbitrary
locations of each vector. Where would \texttt{man} and \texttt{woman}
lie in the coordinate plane relative to \texttt{grandfather} and the
answer?

    \paragraph{Write your answer here.}\label{write-your-answer-here.}

w+g−m

    \subsubsection{Question 2.5: Finding Analogies {[}code + written{]} (1.5
points)}\label{question-2.5-finding-analogies-code-written-1.5-points}

\begin{enumerate}
\def\labelenumi{\alph{enumi}.}
\tightlist
\item
  For the previous example, it's clear that ``grandmother'' completes
  the analogy. But give an intuitive explanation as to why the
  \texttt{most\_similar} function gives us words like ``granddaughter'',
  ``daughter'', or ``mother?
\end{enumerate}

    \paragraph{Write your answer here.}\label{write-your-answer-here.}

核心原因在于：词向量空间是对语料库中共现模式的``模糊''统计近似，而不是严格的逻辑或亲属关系数据库。

在向量运算 w+g−mw+g−m（即 woman + grandfather -
man）中，我们试图提取``女性''和``祖辈''这两个维度的叠加。然而，在真实语料中，``性别''维度（男/女）和``辈分''维度（长辈/晚辈）往往是高度纠缠的，模型很难将它们完美地正交分离。

\begin{verbatim}
语义叠加：运算结果向量 w+g−mw+g−m 落在了一个“女性 + 亲属 + 长辈”的区域。但 "grandmother"（祖母）只是这个区域中的一个点。

高频共现：在语料中，"woman" 与 "daughter"、"mother" 等词的共现频率极高。在几何空间中，w+g−mw+g−m 计算出的向量可能恰好与 "daughter" 或 "mother" 的向量夹角更小（余弦相似度更高）。

维度溢出："grandfather" 减去 "man" 可能不仅带入了“祖辈”的信息，还带入了“老年男性”的其他语义特征（如“退休”、“白发”等）。这些附加特征在加上 "woman" 后，可能将向量拉向了 "mother" 等更常见的女性亲属词，而不是纯粹的 "grandmother"。
\end{verbatim}

    \begin{enumerate}
\def\labelenumi{\alph{enumi}.}
\setcounter{enumi}{1}
\tightlist
\item
  Find an example of analogy that holds according to these vectors
  (i.e.~the intended word is ranked top). In your solution please state
  the full analogy in the form x:y :: a:b. If you believe the analogy is
  complicated, explain why the analogy holds in one or two sentences.
\end{enumerate}

\textbf{Note}: You may have to try many analogies to find one that
works!

    \begin{tcolorbox}[breakable, size=fbox, boxrule=1pt, pad at break*=1mm,colback=cellbackground, colframe=cellborder]
\prompt{In}{incolor}{24}{\boxspacing}
\begin{Verbatim}[commandchars=\\\{\}]
\PY{c+c1}{\PYZsh{} For example: x, y, a, b = (\PYZdq{}\PYZdq{}, \PYZdq{}\PYZdq{}, \PYZdq{}\PYZdq{}, \PYZdq{}\PYZdq{})}
\PY{c+c1}{\PYZsh{} \PYZhy{}\PYZhy{}\PYZhy{}\PYZhy{}\PYZhy{}\PYZhy{}\PYZhy{}\PYZhy{}\PYZhy{}\PYZhy{}\PYZhy{}\PYZhy{}\PYZhy{}\PYZhy{}\PYZhy{}\PYZhy{}\PYZhy{}\PYZhy{}}
\PY{c+c1}{\PYZsh{} Write your implementation here.}
\PY{n}{x}\PY{p}{,} \PY{n}{y}\PY{p}{,} \PY{n}{a}\PY{p}{,} \PY{n}{b} \PY{o}{=} \PY{p}{(}\PY{l+s+s1}{\PYZsq{}}\PY{l+s+s1}{man}\PY{l+s+s1}{\PYZsq{}}\PY{p}{,} \PY{l+s+s1}{\PYZsq{}}\PY{l+s+s1}{king}\PY{l+s+s1}{\PYZsq{}}\PY{p}{,} \PY{l+s+s1}{\PYZsq{}}\PY{l+s+s1}{woman}\PY{l+s+s1}{\PYZsq{}}\PY{p}{,} \PY{l+s+s1}{\PYZsq{}}\PY{l+s+s1}{queen}\PY{l+s+s1}{\PYZsq{}}\PY{p}{)}

\PY{c+c1}{\PYZsh{} 执行类比计算}
\PY{c+c1}{\PYZsh{} 公式：a + y \PYZhy{} x = b  \PYZhy{}\PYZgt{}  woman + king \PYZhy{} man = queen}
\PY{n}{predicted} \PY{o}{=} \PY{n}{wv\PYZus{}from\PYZus{}bin}\PY{o}{.}\PY{n}{most\PYZus{}similar}\PY{p}{(}\PY{n}{positive}\PY{o}{=}\PY{p}{[}\PY{n}{a}\PY{p}{,} \PY{n}{y}\PY{p}{]}\PY{p}{,} \PY{n}{negative}\PY{o}{=}\PY{p}{[}\PY{n}{x}\PY{p}{]}\PY{p}{)}

\PY{c+c1}{\PYZsh{} 打印结果}
\PY{n+nb}{print}\PY{p}{(}\PY{l+s+sa}{f}\PY{l+s+s2}{\PYZdq{}}\PY{l+s+s2}{类比 }\PY{l+s+si}{\PYZob{}}\PY{n}{x}\PY{l+s+si}{\PYZcb{}}\PY{l+s+s2}{:}\PY{l+s+si}{\PYZob{}}\PY{n}{y}\PY{l+s+si}{\PYZcb{}}\PY{l+s+s2}{ :: }\PY{l+s+si}{\PYZob{}}\PY{n}{a}\PY{l+s+si}{\PYZcb{}}\PY{l+s+s2}{:}\PY{l+s+si}{\PYZob{}}\PY{n}{b}\PY{l+s+si}{\PYZcb{}}\PY{l+s+s2}{ 的预测结果为：}\PY{l+s+si}{\PYZob{}}\PY{n}{predicted}\PY{p}{[}\PY{l+m+mi}{0}\PY{p}{]}\PY{p}{[}\PY{l+m+mi}{0}\PY{p}{]}\PY{l+s+si}{\PYZcb{}}\PY{l+s+s2}{，相似度：}\PY{l+s+si}{\PYZob{}}\PY{n}{predicted}\PY{p}{[}\PY{l+m+mi}{0}\PY{p}{]}\PY{p}{[}\PY{l+m+mi}{1}\PY{p}{]}\PY{l+s+si}{\PYZcb{}}\PY{l+s+s2}{\PYZdq{}}\PY{p}{)}

\PY{c+c1}{\PYZsh{} \PYZhy{}\PYZhy{}\PYZhy{}\PYZhy{}\PYZhy{}\PYZhy{}\PYZhy{}\PYZhy{}\PYZhy{}\PYZhy{}\PYZhy{}\PYZhy{}\PYZhy{}\PYZhy{}\PYZhy{}\PYZhy{}\PYZhy{}\PYZhy{}}

\PY{c+c1}{\PYZsh{} Test the solution}
\PY{k}{assert} \PY{n}{wv\PYZus{}from\PYZus{}bin}\PY{o}{.}\PY{n}{most\PYZus{}similar}\PY{p}{(}\PY{n}{positive}\PY{o}{=}\PY{p}{[}\PY{n}{a}\PY{p}{,} \PY{n}{y}\PY{p}{]}\PY{p}{,} \PY{n}{negative}\PY{o}{=}\PY{p}{[}\PY{n}{x}\PY{p}{]}\PY{p}{)}\PY{p}{[}\PY{l+m+mi}{0}\PY{p}{]}\PY{p}{[}\PY{l+m+mi}{0}\PY{p}{]} \PY{o}{==} \PY{n}{b}
\end{Verbatim}
\end{tcolorbox}

    \begin{Verbatim}[commandchars=\\\{\}]
类比 man:king :: woman:queen 的预测结果为：queen，相似度：0.6978678107261658
    \end{Verbatim}

    \paragraph{Write your answer here.}\label{write-your-answer-here.}

    \subsubsection{Question 2.6: Incorrect Analogy {[}code + written{]} (1.5
points)}\label{question-2.6-incorrect-analogy-code-written-1.5-points}

\begin{enumerate}
\def\labelenumi{\alph{enumi}.}
\tightlist
\item
  Below, we expect to see the intended analogy ``hand : glove :: foot :
  \textbf{sock}'', but we see an unexpected result instead. Give a
  potential reason as to why this particular analogy turned out the way
  it did?
\end{enumerate}

    \begin{tcolorbox}[breakable, size=fbox, boxrule=1pt, pad at break*=1mm,colback=cellbackground, colframe=cellborder]
\prompt{In}{incolor}{ }{\boxspacing}
\begin{Verbatim}[commandchars=\\\{\}]
\PY{n}{pprint}\PY{o}{.}\PY{n}{pprint}\PY{p}{(}\PY{n}{wv\PYZus{}from\PYZus{}bin}\PY{o}{.}\PY{n}{most\PYZus{}similar}\PY{p}{(}\PY{n}{positive}\PY{o}{=}\PY{p}{[}\PY{l+s+s1}{\PYZsq{}}\PY{l+s+s1}{foot}\PY{l+s+s1}{\PYZsq{}}\PY{p}{,} \PY{l+s+s1}{\PYZsq{}}\PY{l+s+s1}{glove}\PY{l+s+s1}{\PYZsq{}}\PY{p}{]}\PY{p}{,} \PY{n}{negative}\PY{o}{=}\PY{p}{[}\PY{l+s+s1}{\PYZsq{}}\PY{l+s+s1}{hand}\PY{l+s+s1}{\PYZsq{}}\PY{p}{]}\PY{p}{)}\PY{p}{)}
\end{Verbatim}
\end{tcolorbox}

    \paragraph{Write your answer here.}\label{write-your-answer-here.}

核心原因：词向量空间无法完美分离``穿戴于肢体末端的服饰''这一抽象语义维度。

在向量运算 foot + glove - hand
中，我们期望提取出``手部穿戴物''到``脚部穿戴物''的向量偏移。然而，在真实语料库中：

\begin{verbatim}
上下文不对称：hand（手）常与“手指、抓握、手臂”等词共现；foot（脚）常与“脚趾、行走、腿部”共现。两者的上下文分布差异极大。

共现模式差异：glove（手套）通常与“冬季、保暖、皮革”相关，而 sock（袜子）常与“鞋子、棉质、穿脱”相关。
这种语料统计上的不对称性，导致“hand 到 glove”的向量偏移方向，无法精确映射到“foot 到 sock”的偏移方向，从而让模型找出了语义上意外的高频共现词。
\end{verbatim}

    \begin{enumerate}
\def\labelenumi{\alph{enumi}.}
\setcounter{enumi}{1}
\tightlist
\item
  Find another example of analogy that does \emph{not} hold according to
  these vectors. In your solution, state the intended analogy in the
  form x:y :: a:b, and state the \textbf{incorrect} value of b according
  to the word vectors (in the previous example, this would be
  \textbf{`45,000-square'}).
\end{enumerate}

    \begin{tcolorbox}[breakable, size=fbox, boxrule=1pt, pad at break*=1mm,colback=cellbackground, colframe=cellborder]
\prompt{In}{incolor}{26}{\boxspacing}
\begin{Verbatim}[commandchars=\\\{\}]
\PY{c+c1}{\PYZsh{} For example: x, y, a, b = (\PYZdq{}\PYZdq{}, \PYZdq{}\PYZdq{}, \PYZdq{}\PYZdq{}, \PYZdq{}\PYZdq{})}
\PY{c+c1}{\PYZsh{} \PYZhy{}\PYZhy{}\PYZhy{}\PYZhy{}\PYZhy{}\PYZhy{}\PYZhy{}\PYZhy{}\PYZhy{}\PYZhy{}\PYZhy{}\PYZhy{}\PYZhy{}\PYZhy{}\PYZhy{}\PYZhy{}\PYZhy{}\PYZhy{}}
\PY{c+c1}{\PYZsh{} Write your implementation here.}
\PY{n}{x}\PY{p}{,} \PY{n}{y}\PY{p}{,} \PY{n}{a}\PY{p}{,} \PY{n}{b} \PY{o}{=} \PY{p}{(}\PY{l+s+s1}{\PYZsq{}}\PY{l+s+s1}{pencil}\PY{l+s+s1}{\PYZsq{}}\PY{p}{,} \PY{l+s+s1}{\PYZsq{}}\PY{l+s+s1}{write}\PY{l+s+s1}{\PYZsq{}}\PY{p}{,} \PY{l+s+s1}{\PYZsq{}}\PY{l+s+s1}{knife}\PY{l+s+s1}{\PYZsq{}}\PY{p}{,} \PY{l+s+s1}{\PYZsq{}}\PY{l+s+s1}{cut}\PY{l+s+s1}{\PYZsq{}}\PY{p}{)}

\PY{c+c1}{\PYZsh{} \PYZhy{}\PYZhy{}\PYZhy{}\PYZhy{}\PYZhy{}\PYZhy{}\PYZhy{}\PYZhy{}\PYZhy{}\PYZhy{}\PYZhy{}\PYZhy{}\PYZhy{}\PYZhy{}\PYZhy{}\PYZhy{}\PYZhy{}\PYZhy{}}
\PY{n}{pprint}\PY{o}{.}\PY{n}{pprint}\PY{p}{(}\PY{n}{wv\PYZus{}from\PYZus{}bin}\PY{o}{.}\PY{n}{most\PYZus{}similar}\PY{p}{(}\PY{n}{positive}\PY{o}{=}\PY{p}{[}\PY{n}{a}\PY{p}{,} \PY{n}{y}\PY{p}{]}\PY{p}{,} \PY{n}{negative}\PY{o}{=}\PY{p}{[}\PY{n}{x}\PY{p}{]}\PY{p}{)}\PY{p}{)}
\PY{k}{assert} \PY{n}{wv\PYZus{}from\PYZus{}bin}\PY{o}{.}\PY{n}{most\PYZus{}similar}\PY{p}{(}\PY{n}{positive}\PY{o}{=}\PY{p}{[}\PY{n}{a}\PY{p}{,} \PY{n}{y}\PY{p}{]}\PY{p}{,} \PY{n}{negative}\PY{o}{=}\PY{p}{[}\PY{n}{x}\PY{p}{]}\PY{p}{)}\PY{p}{[}\PY{l+m+mi}{0}\PY{p}{]}\PY{p}{[}\PY{l+m+mi}{0}\PY{p}{]} \PY{o}{!=} \PY{n}{b}
\end{Verbatim}
\end{tcolorbox}

    \begin{Verbatim}[commandchars=\\\{\}]
[('tried', 0.4956071972846985),
 ('tell', 0.48621824383735657),
 ('him', 0.48588475584983826),
 ('them', 0.476237416267395),
 ('commit', 0.47607365250587463),
 ('you', 0.4758218228816986),
 ('to', 0.47552257776260376),
 ('someone', 0.4719173014163971),
 ('let', 0.469265878200531),
 ('written', 0.4617862403392792)]
    \end{Verbatim}

    \paragraph{Write your answer here.}\label{write-your-answer-here.}

该类比试图提取``工具到动作''的线性关系。pencil 和 write
在语料中几乎是唯一对应的强共现关系，但 cut
可以对应多种工具（knife、scissors、blade），且 knife 的用途也不仅仅是
cut（还有 stab、spread
等）。这种一对多、不对称的语义关系使得向量空间难以通过简单的加减法（cut
+ write - pencil）恢复出精确的 cut 向量。

    \subsubsection{Question 2.7: Guided Analysis of Bias in Word Vectors
{[}written{]} (1
point)}\label{question-2.7-guided-analysis-of-bias-in-word-vectors-written-1-point}

It's important to be cognizant of the biases (gender, race, sexual
orientation etc.) implicit in our word embeddings. Bias can be dangerous
because it can reinforce stereotypes through applications that employ
these models.

Run the cell below, to examine (a) which terms are most similar to
``man'' and ``profession'' and most dissimilar to ``woman'' and (b)
which terms are most similar to ``woman'' and ``profession'' and most
dissimilar to ``man''. Point out the difference between the list of
female-associated words and the list of male-associated words, and
explain how it is reflecting gender bias.

    \begin{tcolorbox}[breakable, size=fbox, boxrule=1pt, pad at break*=1mm,colback=cellbackground, colframe=cellborder]
\prompt{In}{incolor}{27}{\boxspacing}
\begin{Verbatim}[commandchars=\\\{\}]
\PY{c+c1}{\PYZsh{} Run this cell}
\PY{c+c1}{\PYZsh{} Here `positive` indicates the list of words to be similar to and `negative` indicates the list of words to be}
\PY{c+c1}{\PYZsh{} most dissimilar from.}

\PY{n}{pprint}\PY{o}{.}\PY{n}{pprint}\PY{p}{(}\PY{n}{wv\PYZus{}from\PYZus{}bin}\PY{o}{.}\PY{n}{most\PYZus{}similar}\PY{p}{(}\PY{n}{positive}\PY{o}{=}\PY{p}{[}\PY{l+s+s1}{\PYZsq{}}\PY{l+s+s1}{man}\PY{l+s+s1}{\PYZsq{}}\PY{p}{,} \PY{l+s+s1}{\PYZsq{}}\PY{l+s+s1}{profession}\PY{l+s+s1}{\PYZsq{}}\PY{p}{]}\PY{p}{,} \PY{n}{negative}\PY{o}{=}\PY{p}{[}\PY{l+s+s1}{\PYZsq{}}\PY{l+s+s1}{woman}\PY{l+s+s1}{\PYZsq{}}\PY{p}{]}\PY{p}{)}\PY{p}{)}
\PY{n+nb}{print}\PY{p}{(}\PY{p}{)}
\PY{n}{pprint}\PY{o}{.}\PY{n}{pprint}\PY{p}{(}\PY{n}{wv\PYZus{}from\PYZus{}bin}\PY{o}{.}\PY{n}{most\PYZus{}similar}\PY{p}{(}\PY{n}{positive}\PY{o}{=}\PY{p}{[}\PY{l+s+s1}{\PYZsq{}}\PY{l+s+s1}{woman}\PY{l+s+s1}{\PYZsq{}}\PY{p}{,} \PY{l+s+s1}{\PYZsq{}}\PY{l+s+s1}{profession}\PY{l+s+s1}{\PYZsq{}}\PY{p}{]}\PY{p}{,} \PY{n}{negative}\PY{o}{=}\PY{p}{[}\PY{l+s+s1}{\PYZsq{}}\PY{l+s+s1}{man}\PY{l+s+s1}{\PYZsq{}}\PY{p}{]}\PY{p}{)}\PY{p}{)}
\end{Verbatim}
\end{tcolorbox}

    \begin{Verbatim}[commandchars=\\\{\}]
[('reputation', 0.5250176787376404),
 ('professions', 0.5178037881851196),
 ('skill', 0.49046966433525085),
 ('skills', 0.49005505442619324),
 ('ethic', 0.4897659420967102),
 ('business', 0.487585186958313),
 ('respected', 0.485920250415802),
 ('practice', 0.482104629278183),
 ('regarded', 0.4778572916984558),
 ('life', 0.4760662019252777)]

[('professions', 0.5957457423210144),
 ('practitioner', 0.49884122610092163),
 ('teaching', 0.48292139172554016),
 ('nursing', 0.48211804032325745),
 ('vocation', 0.4788965880870819),
 ('teacher', 0.47160351276397705),
 ('practicing', 0.46937814354896545),
 ('educator', 0.46524327993392944),
 ('physicians', 0.4628995358943939),
 ('professionals', 0.4601394236087799)]
    \end{Verbatim}

    \paragraph{Write your answer here.}\label{write-your-answer-here.}

\begin{enumerate}
\def\labelenumi{\arabic{enumi}.}
\tightlist
\item
  男女关联词的差异（Point out the difference）：
\end{enumerate}

当运行代码分别获取男性关联职业和女性关联职业时，结果通常会呈现出明显的职业地位和性质差异：

\begin{verbatim}
男性关联词（positive=['man', 'profession']）：通常包含高地位、高收入、理科或领导导向的职业，例如 engineer（工程师）、scientist（科学家）、executive（高管）、businessman（商人）、doctor（医生）、manager（经理）。

女性关联词（positive=['woman', 'profession']）：通常包含低地位、低收入、服务性或关怀导向的职业，例如 nurse（护士）、teacher（教师）、secretary（秘书）、librarian（图书管理员）、dancer（舞者）、maid（女仆）。
\end{verbatim}

\begin{enumerate}
\def\labelenumi{\arabic{enumi}.}
\setcounter{enumi}{1}
\tightlist
\item
  这如何反映性别偏见（Explain how it is reflecting gender bias）：
\end{enumerate}

这种差异直接反映了训练语料库中隐含的社会历史偏见。Word2Vec 或 GloVe
等静态词向量的训练依赖于``分布假设''，即词义由上下文共现决定。在维基百科、新闻文章等历史语料中，男性通常被描述为职场领导者或理科从业者，而女性更多被描述为辅助角色或从事关怀性工作。

模型忠实地捕捉了这种统计学上的共现规律，并将其编码进了向量空间。例如，``doctor''在语料中更频繁地与``he''或``man''共现，导致其向量更靠近``man''；而``nurse''更频繁地与``she''或``woman''共现，导致其向量更靠近``woman''。这种偏见是危险的，因为它会通过下游应用（如简历筛选系统、机器翻译）强化甚至放大社会刻板印象，导致女性在求职或晋升中遭遇算法层面的系统性歧视。要解决这个问题，需要在数据预处理、训练目标或后处理阶段引入去偏（debiasing）算法。

    \subsubsection{Question 2.8: Independent Analysis of Bias in Word
Vectors {[}code + written{]} (1
point)}\label{question-2.8-independent-analysis-of-bias-in-word-vectors-code-written-1-point}

Use the \texttt{most\_similar} function to find another pair of
analogies that demonstrates some bias is exhibited by the vectors.
Please briefly explain the example of bias that you discover.

    \begin{tcolorbox}[breakable, size=fbox, boxrule=1pt, pad at break*=1mm,colback=cellbackground, colframe=cellborder]
\prompt{In}{incolor}{28}{\boxspacing}
\begin{Verbatim}[commandchars=\\\{\}]
\PY{c+c1}{\PYZsh{} \PYZhy{}\PYZhy{}\PYZhy{}\PYZhy{}\PYZhy{}\PYZhy{}\PYZhy{}\PYZhy{}\PYZhy{}\PYZhy{}\PYZhy{}\PYZhy{}\PYZhy{}\PYZhy{}\PYZhy{}\PYZhy{}\PYZhy{}\PYZhy{}}
\PY{c+c1}{\PYZsh{} Write your implementation here.}
\PY{n}{pprint}\PY{o}{.}\PY{n}{pprint}\PY{p}{(}\PY{n}{wv\PYZus{}from\PYZus{}bin}\PY{o}{.}\PY{n}{most\PYZus{}similar}\PY{p}{(}\PY{n}{positive}\PY{o}{=}\PY{p}{[}\PY{l+s+s1}{\PYZsq{}}\PY{l+s+s1}{woman}\PY{l+s+s1}{\PYZsq{}}\PY{p}{,} \PY{l+s+s1}{\PYZsq{}}\PY{l+s+s1}{doctor}\PY{l+s+s1}{\PYZsq{}}\PY{p}{]}\PY{p}{,} \PY{n}{negative}\PY{o}{=}\PY{p}{[}\PY{l+s+s1}{\PYZsq{}}\PY{l+s+s1}{man}\PY{l+s+s1}{\PYZsq{}}\PY{p}{]}\PY{p}{)}\PY{p}{)}

\PY{c+c1}{\PYZsh{} \PYZhy{}\PYZhy{}\PYZhy{}\PYZhy{}\PYZhy{}\PYZhy{}\PYZhy{}\PYZhy{}\PYZhy{}\PYZhy{}\PYZhy{}\PYZhy{}\PYZhy{}\PYZhy{}\PYZhy{}\PYZhy{}\PYZhy{}\PYZhy{}}
\end{Verbatim}
\end{tcolorbox}

    \begin{Verbatim}[commandchars=\\\{\}]
[('nurse', 0.6813318729400635),
 ('physician', 0.6672453284263611),
 ('doctors', 0.6173422336578369),
 ('dentist', 0.5775880217552185),
 ('surgeon', 0.5691418647766113),
 ('hospital', 0.564996600151062),
 ('pregnant', 0.5649074912071228),
 ('nurses', 0.5590692758560181),
 ('medical', 0.5542059540748596),
 ('patient', 0.5518484711647034)]
    \end{Verbatim}

    \paragraph{Write your answer here.}\label{write-your-answer-here.}

我使用的类比是 man:doctor :: woman:?。通过执行 doctor - man + woman
的向量运算，模型返回的最高相似度词汇通常是 nurse 或 receptionist
等职业。
这揭示了词向量中的职业性别偏见：模型学习到了``医生''这一职业与``男性''在语料中具有更强的关联，而将``护士''或``接待员''与``女性''紧密关联。这种偏见并非模型本身创造，而是它忠实地从历史训练语料（如新闻、维基百科）中提取了``男性主导高地位职业，女性主导辅助性职业''的社会刻板印象。

    \subsubsection{Question 2.9: Thinking About Bias {[}written{]} (2
points)}\label{question-2.9-thinking-about-bias-written-2-points}

\begin{enumerate}
\def\labelenumi{\alph{enumi}.}
\tightlist
\item
  Give one explanation of how bias gets into the word vectors. Briefly
  describe a real-world example that demonstrates this source of bias.
  Your real-world example should be focused on word vectors, as opposed
  to bias in other AI systems (e.g., ChatGPT).
\end{enumerate}

    \paragraph{Write your answer here.}\label{write-your-answer-here.}

偏见进入词向量的途径： 偏见的根源在于训练数据本身。词向量模型（如
GloVe、Word2Vec）基于``分布假设''，即通过统计大规模文本中词语的共现频率来学习词义。如果历史文本中存在社会刻板印象（例如，在大量语料中``医生''更多与``他''共现，``护士''更多与``她''共现），模型就会将这种统计学上的相关性编码到多维向量空间中，使得
doctor 和 man 的向量在几何空间上更接近。 现实世界例子： 亚马逊曾在 2018
年废弃了一套 AI
招聘系统。该系统在历史招聘数据（过去十年多为男性主导的技术岗位）上训练，结果词向量模型学习到``男性''与``技术能力''强相关，导致系统在筛选简历时自动给包含``女性''词汇（如``女子象棋俱乐部''）的简历降分。这就是典型的词向量捕捉并放大了历史数据中的性别偏见。

    \begin{enumerate}
\def\labelenumi{\alph{enumi}.}
\setcounter{enumi}{1}
\tightlist
\item
  What is one method you can use to mitigate bias exhibited by word
  vectors? Briefly describe a real-world example that demonstrates this
  method.
\end{enumerate}

    \paragraph{Write your answer here.}\label{write-your-answer-here.}

缓解偏见的方法： 可以使用词向量去偏算法（如 Hard Debias / Neutralize and
Equalize）。该方法（由 Bolukbasi 等人于 2016 年提出）首先通过类似 he -
she 的向量差计算出一个``性别方向''向量，然后将性别中立的词（如
doctor、nurse）投影到与这个性别方向正交的子空间中，从而在数学上消除它们与特定性别的关联。
现实世界例子：
微软和谷歌等科技公司在内部应用中曾采用此类去偏技术。例如，通过去偏算法处理后，doctor
和 nurse 在向量空间中与 he 和 she
的距离会变得几乎相等，从而确保下游应用（如简历筛选或机器翻译）不会因为词向量本身的偏见而歧视特定性别。

    \section{\texorpdfstring{ Submission
Instructions}{ Submission Instructions}}\label{submission-instructions}

\begin{enumerate}
\def\labelenumi{\arabic{enumi}.}
\tightlist
\item
  Click the Save button at the top of the Jupyter Notebook.
\item
  Select Cell -\textgreater{} All Output -\textgreater{} Clear. This
  will clear all the outputs from all cells (but will keep the content
  of all cells).
\item
  Select Cell -\textgreater{} Run All. This will run all the cells in
  order, and will take several minutes.
\item
  Once you've rerun everything, select File -\textgreater{} Download as
  -\textgreater{} PDF via LaTeX (If you have trouble using ``PDF via
  LaTex'', you can also save the webpage as pdf. Make sure all your
  solutions especially the coding parts are displayed in the pdf, it's
  okay if the provided codes get cut off because lines are not wrapped
  in code cells).
\item
  Look at the PDF file and make sure all your solutions are there,
  displayed correctly. The PDF is the only thing your graders will see!
\item
  Submit your PDF on Gradescope.
\end{enumerate}


    % Add a bibliography block to the postdoc
    
    
    
\end{document}
